% ═══════════════════════════════════════════════════════════════════════════
% Master-thesis chapter: Signal Processing and Data Analysis
% ---------------------------------------------------------------------------
% Self-contained \section-level source; include from the thesis main file with
%   % ═══════════════════════════════════════════════════════════════════════════
% Master-thesis chapter: Signal Processing and Data Analysis
% ---------------------------------------------------------------------------
% Self-contained \section-level source; include from the thesis main file with
%   % ═══════════════════════════════════════════════════════════════════════════
% Master-thesis chapter: Signal Processing and Data Analysis
% ---------------------------------------------------------------------------
% Self-contained \section-level source; include from the thesis main file with
%   % ═══════════════════════════════════════════════════════════════════════════
% Master-thesis chapter: Signal Processing and Data Analysis
% ---------------------------------------------------------------------------
% Self-contained \section-level source; include from the thesis main file with
%   \input{processing_chapter}
% Required packages: amsmath, amssymb, siunitx, booktabs (for the tables).
% Every equation cites the implementing function in ldv_analysis/ in a margin
% comment so the text can be checked against the code.
% ═══════════════════════════════════════════════════════════════════════════

\section{Signal Processing and Data Analysis}
\label{sec:processing}

All quantitative results in this thesis are derived from the raw Laser
Doppler Vibrometer (LDV) recordings through the processing chain described in
this chapter. The chain is implemented in the Python package
\texttt{ldv\_analysis}; each subsection below corresponds to one module of
that package, and equations are stated exactly as implemented.

\subsection{Measurement Data and Preprocessing}
\label{sec:proc:preprocessing}

Each measurement consists of two synchronous recordings: the scanning LDV
measures the three-component velocity field
$\mathbf{v}(x_i, t) = (v_x, v_y, v_z)$ at $N_s$ scan points along the cable,
and a reference channel records the three-component velocity of the shaker
head, $\mathbf{v}_\mathrm{sh}(t)$. The excitation is a logarithmic sweep from
$f_1 = \SI{1}{Hz}$ to $f_2 = \SI{500}{Hz}$ over $T = \SI{10}{s}$ (after a
\SI{1}{s} buffer), sampled at $f_s = \SI{5}{kHz}$ with $N = 60\,000$ samples.
The instantaneous sweep frequency reaches $f$ at
% ldv_analysis/config.py :: sweep_time_of_frequency
\begin{equation}
  t_c(f) \;=\; t_\mathrm{buf} + T\,
    \frac{\ln\!\left(f/f_1\right)}{\ln\!\left(f_2/f_1\right)},
  \label{eq:sweeptime}
\end{equation}
which is used to associate a time window with each excitation frequency.

Preprocessing consists of sorting the scan points along the cable axis,
removing dead channels (all-zero velocity), selecting the active shaker
channel by maximum energy, and correcting the known offset between the LDV
and shaker coordinate systems.

\paragraph{Band-pass filtering.}
% ldv_analysis/signal.py :: bandpass
Where a single excitation frequency $f_0$ is analysed, the velocities are
band-pass filtered with a zero-phase Butterworth filter of order 4
(applied forward and backward, \texttt{sosfiltfilt}, so the effective order
is 8 and the phase response is identically zero). The passband is
$f_0 \pm \SI{5}{Hz}$ (absolute bandwidth mode). Zero-phase filtering is
essential here: the coupling efficiency is a ratio of two simultaneously
filtered signals, and any filter-induced phase shift would bias it.

\paragraph{Integration to displacement.}
% ldv_analysis/signal.py :: integrate_fft
The LDV measures velocity; all coupling quantities are defined on
displacements. Integration is performed in the frequency domain. With
$V(f) = \mathcal{F}\{v\}(f)$ the discrete Fourier transform of the velocity,
the displacement spectrum follows from the differentiation property of the
Fourier transform,
\begin{equation}
  u(t) \;=\; \mathcal{F}^{-1}\!\left\{ \frac{V(f)}{i\,2\pi f} \right\}(t),
  \qquad f \ge f_\mathrm{min},
  \label{eq:integration}
\end{equation}
where the division is protected below $f_\mathrm{min} = \SI{0.5}{Hz}$
(frequencies below the floor are divided by $i\,2\pi f_\mathrm{min}$ instead)
to prevent the $1/f$ amplification of DC offset and low-frequency drift that
makes naive time-domain integration unstable. This is a linear, phase-exact
operation: a velocity component $v = \hat v\cos(2\pi f t)$ maps to
$u = \tfrac{\hat v}{2\pi f}\sin(2\pi f t)$, as expected.

\subsection{Geometry, Sag Determination and the Coupling Parameter $\Theta$}
\label{sec:proc:geometry}

\paragraph{Chord and axial projection.}
% ldv_analysis/geometry.py :: compute_chord, project_onto_chord
The cable spans a gap between two mounting points; the first and last scan
point inside the gap, $\mathbf{P}_1$ and $\mathbf{P}_2$, define the
\emph{chord}
\begin{equation}
  \hat{\mathbf{e}} \;=\; \frac{\mathbf{P}_2-\mathbf{P}_1}
                              {\lVert\mathbf{P}_2-\mathbf{P}_1\rVert},
  \qquad
  L \;=\; \lVert\mathbf{P}_2-\mathbf{P}_1\rVert .
  \label{eq:chord}
\end{equation}
Axial (in-line) displacement components are obtained by projection onto the
chord,
\begin{equation}
  u_\parallel(x_i, t) \;=\; \mathbf{u}(x_i,t)\cdot\hat{\mathbf{e}},
  \label{eq:projection}
\end{equation}
which is the displacement component a DAS fibre would sense. Projection onto
the chord (rather than taking the $x$-component) avoids mixing transverse
motion into the axial estimate when the cable is not perfectly aligned with
the coordinate axes.

\paragraph{Sag measurement.}
% ldv_analysis/geometry.py :: measure_sag (use_3d=True, use_parabola=True)
For every scan point in the gap, the perpendicular distance from the chord
line is computed in full 3-D:
\begin{equation}
  s_i = (\mathbf{P}_i - \mathbf{P}_1)\cdot\hat{\mathbf{e}},
  \qquad
  d_i = \bigl\lVert (\mathbf{P}_i-\mathbf{P}_1) - s_i\,\hat{\mathbf{e}}
        \bigr\rVert ,
  \label{eq:perpdist}
\end{equation}
where $s_i$ is the arc coordinate along the chord and $d_i$ the perpendicular
offset. Since the endpoints lie on the chord by construction
($d(0) = d(L) = 0$), a parabola \emph{constrained through both endpoints} is
fitted,
\begin{equation}
  d(s) \;\approx\; -a\, s\,(s - L),
  \qquad
  a \;=\; \frac{\sum_i \phi_i\, d_i}{\sum_i \phi_i^2},
  \quad \phi_i = s_i (s_i - L),
  \label{eq:parabolafit}
\end{equation}
a one-parameter linear least-squares problem. The midpoint sag is the vertex
value of the fitted parabola,
\begin{equation}
  w_0 \;=\; d(L/2) \;=\; \frac{a L^2}{4} \Big|_{a<0} ,
  \label{eq:sagvalue}
\end{equation}
and the root-mean-square residual of the fit,
$\sigma_{w_0} = \sqrt{\tfrac{1}{N}\sum_i (d_i - d(s_i))^2}$, is retained as
the sag measurement uncertainty (it captures outlier scan points and
deviations of the true shape from the parabolic approximation). The parabola
is the small-sag shape of a beam/cable under uniform load, so this fit is
consistent with the analytical model of
Section~\ref{sec:proc:geometry:theta}. If the fit yields $a \ge 0$ (no
physical downward arch), the maximum raw perpendicular distance is used
instead.

\paragraph{The coupling parameter $\Theta$.}
\label{sec:proc:geometry:theta}
% ldv_analysis/config.py :: theta_from_sag, theta_from_material
The bending-stress-relief model \citep{ProbstSimone} predicts that the
static strain-transfer efficiency of a suspended segment is
\begin{equation}
  \eta_\mathrm{stat} \;=\; \frac{\delta x_\ell}{\delta L}
                    \;=\; \frac{1}{1+\Theta},
  \qquad
  \Theta \;=\; \frac{EA}{L}\,\frac{(C_1 x_{f0})^2}{E I C_2},
  \label{eq:etatheta}
\end{equation}
where $x_{f0}$ is the gravity-induced midpoint sag,
$C_1 = \int_0^L (\phi')^2\,\mathrm{d}x$ and
$C_2 = \int_0^L (\phi'')^2\,\mathrm{d}x$ are the mode-shape constants of the
first symmetric clamped--clamped deflection shape
$\phi(x) = \tfrac{1}{2}\bigl[1-\cos(2\pi x/L)\bigr]$. Substituting the
tension-free equilibrium sag
\begin{equation}
  x_{f0} \;=\; w_0 \;=\; \frac{\rho A g L^4}{4\pi^4 E I}
  \label{eq:sagtheory}
\end{equation}
into Eq.~\eqref{eq:etatheta} gives the \emph{material form}
\begin{equation}
  \Theta_\mathrm{mat} \;=\; \frac{\rho^2 g^2 A^3 L^8}{128\,\pi^8 E^2 I^3},
  \label{eq:thetamaterial}
\end{equation}
while keeping the \emph{measured} sag $w_0$ gives the \emph{sag form}
\begin{equation}
  \Theta_\mathrm{sag} \;=\; \frac{A\, w_0^2}{8 I}
  \;\overset{\text{circular}}{=}\; \frac{1}{2}\left(\frac{w_0}{r}\right)^2 ,
  \label{eq:thetasag}
\end{equation}
where the last equality holds for a circular cross-section with
$A = \pi r^2$ and $I = \pi r^4/4$. For the rectangular ribbon cable
(width $w$, height $h$, bending about the weak axis) $A = wh$ and
$I = w h^3/12$ are used, so Eq.~\eqref{eq:thetasag} correctly generalises
beyond circular cables. The sag form is preferred throughout this work
because it absorbs the actual (pretension-affected) deployment geometry; the
material form is reported for comparison. Young's modulus enters as the mean
of three independent tensile tests per cable; $\Theta_\mathrm{mat}$ is
evaluated once per test value and reported as mean~$\pm$~standard deviation.

\subsection{Strain and Elongation Estimation}
\label{sec:proc:strain}

The coupling efficiency compares the \emph{axial elongation of the cable}
with the \emph{elongation imposed at its boundaries}. Three independent
estimators of the cable elongation are implemented --- M1 per-segment 3-D
arc-length, M2 spatial gradient and M3 Fourier-domain gradient --- and their
agreement is a built-in consistency check of the processing.

\subsubsection{Method 1: Per-Segment 3-D Arc-Length Elongation}
\label{sec:proc:strain:m1}
% ldv_analysis/strain.py :: strain_3d_arclength
Method 1 estimates the elongation of each segment between two adjacent scan
points directly from their 3-D displacement vectors, without projecting the
whole field onto the global chord first. Let
$\mathbf{r}_i = \mathbf{P}_{i+1} - \mathbf{P}_i$ be the static segment
vector, $L_{0,i} = \lVert\mathbf{r}_i\rVert$ its rest length and
$\hat{\mathbf{e}}_i = \mathbf{r}_i / L_{0,i}$ its \emph{local} unit vector.
The deformed segment length is
\begin{equation}
  L_i(t) \;=\; \bigl\lVert \mathbf{r}_i
      + \Delta\mathbf{u}_i(t) \bigr\rVert ,
  \qquad
  \Delta\mathbf{u}_i(t) = \mathbf{u}_{i+1}(t) - \mathbf{u}_i(t).
\end{equation}
Because the dynamic displacements ($\sim\si{\micro\meter}$) are many orders
of magnitude smaller than the segment lengths ($\sim\si{\milli\meter}$), the
square root is linearised by a first-order Taylor expansion:
\begin{equation}
  L_i(t) - L_{0,i}
  \;=\; L_{0,i}\sqrt{1
        + \frac{2\,\mathbf{r}_i\!\cdot\!\Delta\mathbf{u}_i}{L_{0,i}^2}
        + \frac{\lVert\Delta\mathbf{u}_i\rVert^2}{L_{0,i}^2}} - L_{0,i}
  \;\approx\;
  \hat{\mathbf{e}}_i \cdot \Delta\mathbf{u}_i(t)
  \;\equiv\; \delta d_i(t).
  \label{eq:strainM1}
\end{equation}
Geometrically: \emph{the relative displacement of the two points is projected
onto the local segment direction} $\hat{\mathbf{e}}_i$ — only the component
of relative motion along the segment changes its length to first order; the
perpendicular component merely rotates it (a second-order length effect that
the linearisation deliberately drops, consistent with the small-strain
regime). For a sagged cable the $\hat{\mathbf{e}}_i$ follow the local cable
direction, so Method 1 measures the \emph{arc-length} change, which is
exactly the quantity a strained optical fibre integrates. The local strain
estimate of segment $i$ is $\varepsilon_i = \delta d_i / L_{0,i}$, and the
total cable elongation is the telescoping sum over the gap,
\begin{equation}
  \delta x_\ell(t) \;=\; \sum_{i\,\in\,\mathrm{gap}} \delta d_i(t).
  \label{eq:deltaxl}
\end{equation}
Method 1 is the primary estimator in this work: it makes no assumption
about the smoothness of the displacement field between scan points and is
insensitive to the chord-projection error for sagged geometries.

\subsubsection{Method 2: Spatial Gradient of the Axial Displacement}
% ldv_analysis/strain.py :: strain_spatial_gradient
Axial strain is the spatial derivative of the axial displacement along the
cable,
\begin{equation}
  \varepsilon(s,t) \;=\; \frac{\partial\, u_\parallel(s,t)}{\partial s},
  \label{eq:strainM2def}
\end{equation}
evaluated with second-order central finite differences on the non-uniform
scan-point grid. For interior point $i$ with spacings
$\Delta s_- = s_i - s_{i-1}$ and $\Delta s_+ = s_{i+1} - s_i$:
\begin{equation}
  \varepsilon_i \;=\;
    u_{\parallel,i+1}\,\frac{\Delta s_-}{\Delta s_+(\Delta s_+ + \Delta s_-)}
  + u_{\parallel,i}\,\frac{\Delta s_+ - \Delta s_-}{\Delta s_+ \Delta s_-}
  - u_{\parallel,i-1}\,\frac{\Delta s_+}{\Delta s_-(\Delta s_+ + \Delta s_-)} .
  \label{eq:strainM2}
\end{equation}
This reduces to the familiar $(u_{i+1}-u_{i-1})/(2\Delta s)$ for uniform
spacing and is exact for locally quadratic displacement fields. The total
elongation follows by trapezoidal integration over the gap,
$\delta x_\ell^{(2)}(t) = \int_\mathrm{gap} \varepsilon(s,t)\,\mathrm{d}s$.

\subsubsection{Method 3: Fourier-Domain Spatial Gradient}
% ldv_analysis/strain.py :: strain_fourier_gradient
As a cross-check with different numerical error characteristics, the axial
displacement is resampled onto a uniform grid $s_k$ (cubic interpolation)
and differentiated spectrally using the Fourier differentiation theorem,
\begin{equation}
  \varepsilon(s,t) \;=\;
  \mathcal{F}_s^{-1}\bigl\{\, i k\, \mathcal{F}_s\{u_\parallel\}\,\bigr\},
  \label{eq:strainM3}
\end{equation}
with $k$ the spatial angular wavenumber. Spectral differentiation is exact
for band-limited periodic signals; since the cable displacement is not
periodic over the gap, Gibbs-type artefacts can appear near the endpoints,
which is why only gap-interior stations enter the elongation integral. In
practice Method 3 agrees with Methods 1--2 in the smooth interior and serves
to expose discretisation artefacts of the finite-difference estimate.

\subsubsection{Reference Elongations}
% ldv_analysis/freqdomain.py :: build_coupling_signals (x_source)
Three references for the imposed boundary elongation $\delta L(t)$ are
available:
\begin{itemize}
  \item \textbf{Endpoint pair} (default): the span change measured on the
    cable itself,
    $\delta L = u_\parallel(\mathbf{P}_2) - u_\parallel(\mathbf{P}_1)$.
    It shares its endpoints with $\delta x_\ell$, so the efficiency isolates
    bending stress relief and is immune to cross-instrument calibration.
  \item \textbf{Shaker}: the chord-projected shaker-head displacement
    $\delta L = \mathbf{u}_\mathrm{sh}\cdot\hat{\mathbf{e}}$ (sign-corrected
    so that positive means stretching), i.e.\ the imposed ground motion.
  \item \textbf{Shaker-to-end}: the elongation of the sub-span between the
    shaker and the adjacent endpoint sensor, used to diagnose losses at the
    first contact.
\end{itemize}

\subsection{Time-Domain Coupling Efficiency (Quality Control)}
\label{sec:proc:timedomain}
% ldv_analysis/analysis.py :: per_frequency_qc

For inspection at a single excitation frequency $f_0$, the record is windowed
to $t_c(f_0) \pm n_\mathrm{cyc}/f_0$ (Eq.~\ref{eq:sweeptime},
$n_\mathrm{cyc}=5$), band-passed (Section~\ref{sec:proc:preprocessing}),
integrated to displacement, and the instantaneous efficiency is formed:
\begin{equation}
  \eta(t) \;=\; \frac{\delta x_\ell(t)}{\delta L(t)},
  \qquad
  \eta_\mathrm{med} = \operatorname{median}_t\, \eta(t)
  \quad\text{over } |\delta L(t)| > 0.05\,\max|\delta L| .
  \label{eq:etatime}
\end{equation}
The amplitude gate excludes the zero crossings of the reference, where the
ratio is ill-conditioned; the median suppresses the remaining outliers. A
phase-insensitive variant replaces both signals by their analytic-signal
envelopes (Section~\ref{sec:proc:amplitude}),
$\eta_\mathrm{env}(t) = |\mathcal{H}\{\delta x_\ell\}| /
|\mathcal{H}\{\delta L\}|$: comparing $\eta_\mathrm{med}$ and
$\eta_\mathrm{env}$ distinguishes whether a low efficiency is caused by an
amplitude deficit or by a phase lag between cable and reference.

These time-domain values are used for quality control and for illustrating
the mechanism at single frequencies; all summary efficiencies in the results
chapter are computed in the frequency domain
(Section~\ref{sec:proc:freqdomain}), which uses the whole record at once and
is considerably more stable.

\subsection{Frequency-Domain Transfer Functions}
\label{sec:proc:freqdomain}
% ldv_analysis/freqdomain.py

\subsubsection{Signals}
From the full (unfiltered) sweep record, integrated to displacement, three
signals are built:
\begin{align}
  X(t) &= \delta L(t) && \text{input: boundary span change,} \notag\\
  Y_1(t) &= \delta x_\ell(t) && \text{output 1: cable elongation
                                (Eq.~\ref{eq:deltaxl}),} \notag\\
  Y_2(t) &= u_{\perp,\mathrm{mid}}(t) && \text{output 2: midpoint deflection
                                along the static sag direction.} \notag
\end{align}
$H_\eta(f) = Y_1/X$ is the dynamic strain-transfer efficiency;
$H_\mathrm{mid}(f) = Y_2/X$ measures the intensity of the bending mode, whose
resonance structure explains the frequency dependence of $H_\eta$
\citep{ProbstSimone}.

\subsubsection{Welch Spectral Estimation}
% ldv_analysis/freqdomain.py :: welch_spectra
The record is divided into $K$ segments of length
$n_\mathrm{seg} = 2 f_s = 10\,000$ samples (\SI{2}{s}, frequency resolution
$\Delta f = \SI{0.5}{Hz}$) with 50\,\% overlap, each tapered with a Hann
window $w[n]$. With $X_k(f)$, $Y_k(f)$ the windowed DFTs of segment $k$,
the auto- and cross-spectral densities are the segment averages
\begin{equation}
  S_{XX}(f) = \frac{1}{K}\sum_{k=1}^{K} |X_k(f)|^2 , \qquad
  S_{YY}(f) = \frac{1}{K}\sum_{k=1}^{K} |Y_k(f)|^2 , \qquad
  S_{XY}(f) = \frac{1}{K}\sum_{k=1}^{K} X_k^*(f)\, Y_k(f),
  \label{eq:welch}
\end{equation}
(up to the standard window normalisation, which cancels in every ratio used
below). The averaging over $K \approx 11$ segments reduces the variance of
each spectral estimate by $\approx 1/K$ — this is what makes the Welch-based
transfer functions far smoother than any single-FFT estimate. The
\emph{ordinary coherence}
\begin{equation}
  \gamma^2(f) \;=\; \frac{|S_{XY}(f)|^2}{S_{XX}(f)\, S_{YY}(f)}
  \;\in\; [0, 1]
  \label{eq:coherence}
\end{equation}
measures the fraction of output power linearly explained by the input;
$\gamma^2 < 1$ indicates noise, nonlinearity, or unmeasured inputs.

\subsubsection{FRF Estimators: Direct, $H_1$, $H_2$, $H_v$}
\label{sec:proc:estimators}
% ldv_analysis/freqdomain.py :: direct_transfer, estimate_H

\paragraph{Direct estimate.}
The conceptually simplest frequency response estimate is the ratio of the
full-record windowed Fourier transforms,
\begin{equation}
  H_\mathrm{direct}(f) \;=\;
  \frac{\mathcal{F}\{w\,Y\}(f)}{\mathcal{F}\{w\,X\}(f)} ,
  \label{eq:Hdirect}
\end{equation}
evaluated on the dense grid $\Delta f_\mathrm{d} = f_s/N = \SI{0.083}{Hz}$.
Every frequency bin is a single realisation: the estimate is unbiased in the
noise-free limit but has no variance reduction, so amplitude and phase
scatter strongly wherever excitation or response is weak.

Comparing Eq.~\eqref{eq:Hdirect} directly against $H_1$ would conflate two
distinct effects — the estimator algebra and the amount of spectral averaging.
To separate them, the direct spectra are smoothed along the \emph{frequency}
axis (Daniell smoothing) rather than over time segments: with a normalised
kernel $k$ of $m$ bins,
\begin{equation}
  \tilde S_{AB}(f_i) = \sum_{j} k_j\, A^*(f_{i+j})\, B(f_{i+j}),
  \qquad
  H_\mathrm{direct}^{\mathrm{sm}}(f) = \frac{\tilde S_{XY}(f)}{\tilde S_{XX}(f)} .
  \label{eq:Hdirectsmooth}
\end{equation}
Smoothing the products \emph{before} the division makes the estimate
input-power weighted, so bins with weak $|X|$ contribute little instead of
diverging; this also removes the positive bias that the unweighted ratio
carries at low SNR. A kernel with $\sum_j k_j = 1$ averages
$n_\mathrm{eff} = 1/\sum_j k_j^2$ independent bins, which sets both the
variance reduction and the effective resolution bandwidth. The kernel width is
therefore chosen so that $n_\mathrm{eff}\,\Delta f_\mathrm{d}$ equals the
\emph{equivalent noise bandwidth} of the Welch estimate,
\begin{equation}
  B_\mathrm{ENBW} = f_s \frac{\sum_n w[n]^2}{\left(\sum_n w[n]\right)^2}
  = 1.5\,\Delta f = \SI{0.75}{Hz}
  \quad\text{(Hann)},
  \label{eq:enbw}
\end{equation}
i.e.\ $n_\mathrm{eff} \approx 9$ dense bins. Note that
$B_\mathrm{ENBW} > \Delta f$: the Hann taper spreads each spectral line over
1.5 bins, so a Welch estimate does not resolve to its nominal bin spacing, and
matching $B_\mathrm{ENBW}$ rather than $\Delta f$ is what puts the two
estimates on an equal footing. With this choice the smoothed direct estimate
and $H_1$ agree to within a few percent in $|H|$ while retaining the dense
frequency grid; the residual difference is attributable to the estimator
algebra alone. Both the raw and the smoothed direct curves are shown, so the
figure separates the effect of averaging (raw $\to$ smoothed) from the effect
of the estimator (smoothed $\to$ $H_1$).

\paragraph{$H_1$ — noise on the output.}
Assume the measured output contains additive noise uncorrelated with the
input, $Y = H X + N$. Minimising the output residual over the segment
ensemble,
$\min_H \sum_k |Y_k - H X_k|^2$, gives the least-squares solution
\begin{equation}
  H_1(f) \;=\; \frac{S_{XY}(f)}{S_{XX}(f)} .
  \label{eq:H1}
\end{equation}
Since $\mathbb{E}\{X^* N\} = 0$, the cross-spectrum averages the output noise
away and $H_1$ is unbiased under output noise. If the \emph{input} channel is
noisy, $S_{XX}$ is inflated and $|H_1|$ \emph{under}estimates the true FRF.

\paragraph{$H_2$ — noise on the input.}
Assume instead noise on the input, $X = X_0 + M$ with $Y = H X_0$.
Minimising the input-referenced residual
$\min_H \sum_k |X_k - Y_k / H|^2$ yields
\begin{equation}
  H_2(f) \;=\; \frac{S_{YY}(f)}{S_{YX}(f)}
        \;=\; \frac{S_{YY}(f)}{S_{XY}^*(f)} .
  \label{eq:H2}
\end{equation}
$H_2$ is unbiased under input noise but \emph{over}estimates $|H|$ when the
output is noisy. The two estimators bracket the true FRF,
\begin{equation}
  |H_1(f)| \;\le\; |H(f)| \;\le\; |H_2(f)|,
  \qquad
  \frac{|H_1(f)|}{|H_2(f)|} = \gamma^2(f),
  \label{eq:bracket}
\end{equation}
so their separation is itself a coherence display: the estimators agree
exactly where $\gamma^2 = 1$ and diverge where the data are noisy.

\paragraph{$H_v$ — noise on both channels (total least squares).}
When both channels carry comparable noise — realistic here, since input and
output are both derived from LDV velocities through the same processing —
the total-least-squares estimator distributes the residual over both
channels. Writing the spectral sample covariance of the vector $(X, Y)$ per
frequency bin,
\begin{equation}
  \mathbf{C}(f) \;=\;
  \begin{pmatrix} S_{XX} & S_{XY}^* \\ S_{XY} & S_{YY} \end{pmatrix},
\end{equation}
the TLS solution is the direction of the eigenvector belonging to the
\emph{smallest} eigenvalue $\lambda_-$ of $\mathbf{C}$ (the noise
eigenvalue). For the $2\times 2$ Hermitian matrix this eigenvalue is
\begin{equation}
  \lambda_-(f) \;=\; \frac{S_{XX}+S_{YY}}{2}
    - \sqrt{\left(\frac{S_{XX}-S_{YY}}{2}\right)^{\!2} + |S_{XY}|^2},
\end{equation}
and solving the eigenvector relation for the FRF gives the closed form
\begin{equation}
  H_v(f) \;=\; \frac{S_{YY}(f) - \lambda_-(f)}{S_{XY}^*(f)} .
  \label{eq:Hv}
\end{equation}
One verifies directly that in the noise-free case ($\lambda_- = 0$,
$S_{YY} = |H|^2 S_{XX}$, $S_{XY} = H S_{XX}$) Eq.~\eqref{eq:Hv} returns
exactly $H$, and that $|H_1| \le |H_v| \le |H_2|$ holds bin-by-bin. All
three Welch estimators share the same phase,
$\arg H_1 = \arg H_2 = \arg H_v = \arg S_{XY}$, so the estimator choice
affects only the magnitude.

\paragraph{Phase presentation.}
Phases are reported \emph{wrapped} to $(-180^\circ, 180^\circ]$. Unwrapped
phase curves accumulate spurious $\pm 360^\circ$ jumps whenever
low-coherence bins randomise the phase, which makes them diverge without
physical meaning; the wrapped display is immune to this and the physically
relevant transition (the $\approx 180^\circ$ flip across the bending
resonance) remains fully visible.

\subsubsection{Band-Averaged Efficiency and Coherence Gating}
% ldv_analysis/freqdomain.py :: band_mean_eta
The quasi-static efficiency reported per configuration is the coherence-gated
band average
\begin{equation}
  \bar\eta \;=\; \bigl\langle\, |H_\eta(f)| \,\bigr\rangle_{f \in B},
  \qquad
  B = \{ f \in [f_\mathrm{min}, f_\mathrm{max}] :
         \gamma^2_\eta(f) \ge 0.7 \},
  \label{eq:bandeta}
\end{equation}
with $[f_\mathrm{min}, f_\mathrm{max}] = [5, 25]\,\si{Hz}$ well below the
first resonance. The gate removes noise-dominated bins that would otherwise
bias the average; the standard deviation over $B$ is reported as the error
bar.

\subsubsection{Single-Frequency Cross-Check and Resonance Characterisation}
% ldv_analysis/freqdomain.py :: stft_point_transfer, fit_lorentzian
As an independent validation, point estimates
$H(f_0) = \mathcal{F}\{w\,Y\}(f_0)\,/\,\mathcal{F}\{w\,X\}(f_0)$ are computed
from short Hann-tapered windows centred on the sweep arrival $t_c(f_0)$ and
overlaid on the Welch curves.

The fundamental resonance $f_1$ is identified as the global maximum of
$|H_\mathrm{mid}(f)|$; a driven single-degree-of-freedom magnitude response
\begin{equation}
  |H(f)| \;=\; \frac{G\, f_1^2}
   {\sqrt{\left(f_1^2 - f^2\right)^2 + \left(f_1 f / Q\right)^2}}
  \label{eq:sdof}
\end{equation}
is fitted in a window around the peak, yielding the quality factor $Q$ and
damping ratio $\zeta = 1/(2Q)$. The theoretical clamped--clamped fundamental
\begin{equation}
  f_1^\mathrm{th} \;=\; \frac{\alpha_1^2}{2\pi L^2}\sqrt{\frac{EI}{\rho A}},
  \qquad \alpha_1^2 = 4.730^2 \approx 22.37,
  \label{eq:f1theory}
\end{equation}
is quoted for comparison (cf.~Section~\ref{sec:proc:modal}).

\subsection{Amplitude Extraction}
\label{sec:proc:amplitude}
% ldv_analysis/amplitude.py

Beyond ratios, the absolute displacement and strain amplitudes are needed as
reference values for the numerical simulations. Two extraction routes are
implemented and cross-validated.

\paragraph{Time domain: analytic-signal envelope.}
For a windowed, band-passed trace $x(t)$ the analytic signal is
$x_a(t) = x(t) + i\,\mathcal{H}\{x\}(t)$ with $\mathcal{H}$ the Hilbert
transform; its magnitude $|x_a(t)|$ is the instantaneous envelope. The
representative amplitude is the envelope median (robust against the
filter-edge droop at the window boundaries); envelope maximum and
$\sqrt{2}\,\mathrm{RMS}$ variants are available for comparison.

\paragraph{Frequency domain: band-energy read-out.}
The velocities are integrated once per dataset; at each target frequency
$f_0$ the displacement signal is windowed around $t_c(f_0)$, Hann-tapered
and Fourier transformed. Because the log-sweep smears the tone over several
bins (especially at low frequency, where the window spans many seconds), the
amplitude is recovered from the \emph{energy in the analysis band}
$B = [f_0 - 5, f_0 + 5]\,\si{Hz}$ rather than from a single bin. For a tone
of peak amplitude $A$ under a window $w[n]$, Parseval's theorem gives
\begin{equation}
  A \;=\; 2\sqrt{\frac{\sum_{k \in B} |X_k|^2}{N \sum_n w[n]^2}} ,
  \label{eq:bandenergy}
\end{equation}
which is exact for an isolated in-band tone. This route avoids the
per-frequency band-pass filtering of the full record and is roughly two
orders of magnitude faster, enabling dense frequency grids; it agrees with
the envelope route to within a few percent below $\approx\SI{100}{Hz}$.

\subsection{Extension--Compression Asymmetry}
\label{sec:proc:asymmetry}
% ldv_analysis/asymmetry.py

A sagged segment stiffens when its endpoints are pulled apart (the sag is
straightened) and softens when they are pushed together (the sag deepens);
numerical simulations show correspondingly different responses for the two
deformation directions. A perfectly linear system responds
mirror-symmetrically to the two half-cycles of a sinusoidal drive, so any
measurable asymmetry is a signature of this geometric nonlinearity. Three
complementary detectors are evaluated per excitation frequency:

\begin{enumerate}
  \item \textbf{Half-cycle conditional efficiency.} The samples of the
  reference elongation are split by sign — extension
  ($\delta L > 0$, endpoints separating) versus compression
  ($\delta L < 0$) — with a $25\,\%$-of-maximum amplitude gate excluding the
  zero crossings. The efficiencies
  \begin{equation}
    \eta_\mathrm{ext} = \operatorname*{median}_{\delta L(t) > \delta_g}
      \frac{\delta x_\ell(t)}{\delta L(t)},
    \qquad
    \eta_\mathrm{comp} = \operatorname*{median}_{\delta L(t) < -\delta_g}
      \frac{\delta x_\ell(t)}{\delta L(t)}
    \label{eq:etahalfcycle}
  \end{equation}
  coincide for a linear system; $\eta_\mathrm{ext} \ne \eta_\mathrm{comp}$
  quantifies the directional coupling difference.

  \item \textbf{Peak asymmetry.} From the mean positive-peak amplitude
  $A_+$ and mean negative-peak magnitude $A_-$ of the cable elongation
  $\delta x_\ell(t)$:
  \begin{equation}
    \alpha \;=\; \frac{A_+ - A_-}{A_+ + A_-} \;\in\; [-1, 1],
    \label{eq:peakasym}
  \end{equation}
  with $\alpha > 0$ meaning the cable elongates further during extension
  half-cycles than it shortens during compression.

  \item \textbf{Even-harmonic distortion.} Expanding the response in powers
  of the drive, $y = c_1 x + c_2 x^2 + c_3 x^3 + \dots$ with
  $x = A\cos\omega t$, the quadratic term produces
  $c_2 A^2 \cos^2\omega t = \tfrac{c_2 A^2}{2}(1 + \cos 2\omega t)$ — energy
  at \emph{twice} the drive frequency — whereas the cubic (symmetric) term
  appears at $3\omega$. Since a sign-asymmetric stiffness is exactly a
  quadratic term, the ratio
  \begin{equation}
    R_2 \;=\; \frac{|Y(2 f_0)|}{|Y(f_0)|}
    \label{eq:h2ratio}
  \end{equation}
  isolates the asymmetric part of the response independently of the
  reference channel; $R_3 = |Y(3f_0)|/|Y(f_0)|$ is reported alongside to
  attribute nonlinearity to the symmetric (softening/stiffening) branch.
\end{enumerate}

For the linearity data set (Section~\ref{sec:proc:linearity}) the half-cycle
efficiencies are additionally evaluated in bins of the instantaneous drive
amplitude, revealing whether an asymmetry develops progressively with
amplitude (e.g.\ slack forming on the compression half-cycle only).

\subsection{Modal Analysis}
\label{sec:proc:modal}
% ldv_analysis/modal.py

\subsubsection{Detection Fingerprint and Spectral Smoothing}
% modal.spectral_fingerprint
Transverse resonances are detected on a \emph{fingerprint}: per-sensor
spectra of the transverse velocities $v_y, v_z$ over the gap, collapsed
across sensors into the total-power detection signal
\begin{equation}
  P(f) \;=\; \sum_{s\,\in\,\mathrm{gap}}
      \Bigl( |V_y(f; s)|^2 + |V_z(f; s)|^2 \Bigr).
  \label{eq:fingerprint}
\end{equation}
The per-sensor spectra are Welch estimates
($n_\mathrm{seg} = 2 f_s$, $\Delta f = \SI{0.5}{Hz}$, Hann, 50\,\% overlap,
$K \approx 11$ averages) rather than raw single-record FFTs: the
$\mathrm{var} \propto 1/K$ averaging suppresses the bin-to-bin fluctuations
that previously caused the peak picker to fire on noise wiggles. An optional
moving-average (boxcar) smoothing over a configurable bandwidth can be
applied on top. Resonances are then the peaks of $P(f)$ (normalised to its
in-band maximum) exceeding a prominence threshold with a minimum mutual
spacing (\texttt{scipy.signal.find\_peaks}).

\paragraph{Input-bias correction.}
% modal.shaker_input_spectrum, normalize_by_input
The shaker head does not excite purely axially: its own recording shows
frequency-dependent energy in the transverse components. A spectral peak in
the cable response may therefore reflect a peak in the \emph{excitation}
rather than a cable resonance. Two diagnostics address this. First, the
input directionality is quantified by the per-frequency power fractions of
the shaker-head components,
\begin{equation}
  \phi_c(f) \;=\; \frac{|V_c^\mathrm{sh}(f)|^2}
   {\sum_{c' \in \{x,y,z\}} |V_{c'}^\mathrm{sh}(f)|^2},
  \qquad c \in \{x, y, z\}.
  \label{eq:directionality}
\end{equation}
Second, the detection signal can be normalised by the total shaker input
power, $P_\mathrm{norm}(f) = P(f) / P_\mathrm{sh}(f)$ — an FRF-like quantity
in which input coloration cancels: peaks that survive the normalisation are
genuine output-per-input amplifications (resonances), while peaks that
disappear were inherited from the excitation spectrum.

\subsubsection{Complex Mode-Shape Extraction}
% modal.extract_mode_shape
For each detected resonance $f_n$, the transverse velocities are band-passed
to $f_n(1 \pm 0.1)$, integrated to displacement, and the complex Fourier
coefficient at the bin nearest $f_n$ is read per sensor:
\begin{equation}
  \varphi_y(s) = U_y(f_n; s), \qquad \varphi_z(s) = U_z(f_n; s).
  \label{eq:modeshape}
\end{equation}
The complex-valued shapes retain the relative phase between sensors, so
travelling-wave content and sign flips across nodes are preserved.

\subsubsection{Euler--Bernoulli Reference Modes}
% modal.clamped_clamped_mode
The gap is modelled as a doubly-clamped uniform beam. Free vibration of an
Euler--Bernoulli beam obeys $EI\, w'''' = \rho A\, \omega^2 w$; with clamped
boundary conditions $w(0)=w'(0)=w(L)=w'(L)=0$, non-trivial solutions exist
only at the roots $\alpha_n$ of the frequency equation
\begin{equation}
  \cos\alpha \cosh\alpha \;=\; 1,
  \qquad
  \alpha_n = 4.730,\ 7.853,\ 10.996,\ 14.137,\ 17.279,\ \dots
  \label{eq:freqequation}
\end{equation}
with natural frequencies and mode shapes
\begin{equation}
  f_n \;=\; \frac{\alpha_n^2}{2\pi L^2} \sqrt{\frac{EI}{\rho A}},
  \qquad
  \varphi_n(\xi) = \bigl[\cosh\alpha_n\xi - \cos\alpha_n\xi\bigr]
   - \sigma_n \bigl[\sinh\alpha_n\xi - \sin\alpha_n\xi\bigr],
  \label{eq:ebmodes}
\end{equation}
where $\xi = x/L \in [0,1]$ and
$\sigma_n = (\cosh\alpha_n - \cos\alpha_n)/(\sinh\alpha_n - \sin\alpha_n)$.
The frequency ratios $f_n / f_1 = (\alpha_n/\alpha_1)^2 =
1,\ 2.76,\ 5.40,\ 8.93,\ \dots$ provide a material-free consistency check of
the mode ordering.

\subsubsection{Modal Assurance Criterion}
% modal.mac
Each measured shape is assigned to a theoretical mode by the Modal Assurance
Criterion,
\begin{equation}
  \operatorname{MAC}(\boldsymbol{\varphi}_m, \boldsymbol{\varphi}_t)
  \;=\;
  \frac{\bigl|\boldsymbol{\varphi}_m^{\mathsf{H}}
              \boldsymbol{\varphi}_t\bigr|^2}
       {\bigl(\boldsymbol{\varphi}_m^{\mathsf{H}}\boldsymbol{\varphi}_m\bigr)
        \bigl(\boldsymbol{\varphi}_t^{\mathsf{H}}\boldsymbol{\varphi}_t\bigr)}
  \;\in\; [0, 1],
  \label{eq:mac}
\end{equation}
where $\boldsymbol{\varphi}_m$ is the measured complex shape (the dominant
polarization component, sampled at the sensor positions),
$\boldsymbol{\varphi}_t$ the theoretical shape $\varphi_n(\xi_i)$ evaluated
at the same positions, and $^{\mathsf{H}}$ the conjugate transpose. The MAC
is the squared modulus of the normalised inner product — a
\emph{shape-correlation} coefficient: it equals 1 when the vectors are
collinear (identical shape up to any complex scale), and 0 when they are
orthogonal. Because a global complex factor cancels, the measured shape
needs no phase alignment or amplitude scaling before comparison, and a
$180^\circ$ phase flip across a node is handled correctly. The measured
resonance is assigned the mode order with the largest MAC; low best-MAC
values flag shapes that no beam mode explains (e.g.\ rigid-body-like motion
from mount compliance or shaker cross-talk).

\subsubsection{Polarization Analysis}
% modal.polarization_ellipse
At each sensor the pair $(\varphi_y, \varphi_z)$ defines an ellipse traced
by $\Re\{(\varphi_y, \varphi_z)\, e^{i\omega t}\}$ in the $y$--$z$ plane.
Decomposing into counter-rotating circular components
\begin{equation}
  A_+ = \tfrac{1}{2}\bigl(\varphi_y + i\,\varphi_z\bigr),
  \qquad
  A_- = \tfrac{1}{2}\bigl(\varphi_y^* + i\,\varphi_z^*\bigr)
  \quad\Longrightarrow\quad
  \text{major} = |A_+| + |A_-|,\quad
  \text{minor} = \bigl||A_+| - |A_-|\bigr|,
  \label{eq:ellipse}
\end{equation}
with axis ratio $\mathrm{minor}/\mathrm{major}$ (0 = linear, 1 = circular),
orientation $\theta = \tfrac{1}{2}(\arg A_+ + \arg A_-)$ and handedness
$\operatorname{sign}\Im(\varphi_y \varphi_z^*)$. A resonance is classified
linear / elliptical / circular by the amplitude-weighted majority over the
significant sensors.

\subsection{Linearity Analysis}
\label{sec:proc:linearity}
% ldv_analysis/analysis.py :: linearity_qc

A dedicated test signal probes amplitude dependence: eight discrete
frequencies (5--400\,\si{Hz}), each driven for \SI{3.5}{s} with a linear
amplitude ramp $0 \to 1$ followed by a \SI{0.5}{s} fade. Each segment is
extracted in full, band-passed ($f_0 \pm \SI{5}{Hz}$), integrated, and the
efficiency $\eta(t)$ (Eq.~\ref{eq:etatime}) is evaluated against the
\emph{instantaneous drive amplitude}, defined as the analytic-signal envelope
$a(t) = |\mathcal{H}\{\delta L\}(t)|$. Because the ramp sweeps amplitude
monotonically, plotting the binned median of $\eta$ against $a$ directly
tests linearity: a flat curve means amplitude-independent coupling, a slope
or hysteresis loop indicates nonlinearity (e.g.\ frictional slip at the
mounts above a threshold amplitude). A complementary slip indicator divides
the envelope of a single cable-sensor velocity by the envelope of the shaker
reference; for slip-free coupling this ratio stays constant along the ramp.
The frequency-domain cross-check estimates $H_1$
(Eq.~\ref{eq:H1}) on each mono-frequency segment and compares
$|H_1(f_0)|$ with the time-domain median efficiency.

\subsection{Uncertainty Treatment}
\label{sec:proc:uncertainty}

\begin{itemize}
  \item \textbf{Young's modulus:} three tensile tests per cable;
    $\Theta_\mathrm{mat}$ is evaluated per test and reported as mean $\pm$
    standard deviation. Since $\Theta_\mathrm{mat} \propto E^{-2}$, a
    relative modulus error $\epsilon_E$ propagates as
    $2\,\epsilon_E$ into $\Theta$.
  \item \textbf{Gap length:} the mounting-gap reading error is
    $\sigma_L = \SI{0.5}{cm}$. The material form is extremely sensitive,
    $\Theta_\mathrm{mat} \propto L^8 \Rightarrow
    \sigma_\Theta/\Theta = 8\,\sigma_L/L$ (40\,\% at $L=\SI{10}{cm}$!),
    which is the quantitative reason the sag-based
    $\Theta_\mathrm{sag}$ — independent of $L$ — is preferred.
  \item \textbf{Sag:} the RMS residual of the constrained parabola fit
    (Eq.~\ref{eq:parabolafit}) serves as $\sigma_{w_0}$; with
    $\Theta_\mathrm{sag} \propto w_0^2$,
    $\sigma_\Theta/\Theta = 2\,\sigma_{w_0}/w_0$.
  \item \textbf{Efficiency:} the standard deviation of $|H_\eta|$ over the
    coherence-gated quasi-static band (Eq.~\ref{eq:bandeta}); the
    $H_1$/$H_2$ bracket (Eq.~\ref{eq:bracket}) provides an additional
    systematic error envelope.
\end{itemize}

% ─── Bibliography hooks ─────────────────────────────────────────────────────
% \citep{ProbstSimone}  : S. Probst et al., "Unburied DAS on the Moon:
%                         Modeling Ground-to-Cable Strain Transfer" (draft).
% Probst et al. 2026    : Controlled Source DAS Coupling Tests, Earth and
%                         Space Science 13(3), e2025EA004817.
% Welch 1967            : P. Welch, IEEE Trans. Audio Electroacoust. 15(2).
% Allemang 2003         : R. Allemang, "The Modal Assurance Criterion --
%                         Twenty Years of Use and Abuse", Sound & Vibration.
% Bendat & Piersol 2010 : Random Data: Analysis and Measurement Procedures
%                         (H1/H2/Hv estimators, coherence).

% Required packages: amsmath, amssymb, siunitx, booktabs (for the tables).
% Every equation cites the implementing function in ldv_analysis/ in a margin
% comment so the text can be checked against the code.
% ═══════════════════════════════════════════════════════════════════════════

\section{Signal Processing and Data Analysis}
\label{sec:processing}

All quantitative results in this thesis are derived from the raw Laser
Doppler Vibrometer (LDV) recordings through the processing chain described in
this chapter. The chain is implemented in the Python package
\texttt{ldv\_analysis}; each subsection below corresponds to one module of
that package, and equations are stated exactly as implemented.

\subsection{Measurement Data and Preprocessing}
\label{sec:proc:preprocessing}

Each measurement consists of two synchronous recordings: the scanning LDV
measures the three-component velocity field
$\mathbf{v}(x_i, t) = (v_x, v_y, v_z)$ at $N_s$ scan points along the cable,
and a reference channel records the three-component velocity of the shaker
head, $\mathbf{v}_\mathrm{sh}(t)$. The excitation is a logarithmic sweep from
$f_1 = \SI{1}{Hz}$ to $f_2 = \SI{500}{Hz}$ over $T = \SI{10}{s}$ (after a
\SI{1}{s} buffer), sampled at $f_s = \SI{5}{kHz}$ with $N = 60\,000$ samples.
The instantaneous sweep frequency reaches $f$ at
% ldv_analysis/config.py :: sweep_time_of_frequency
\begin{equation}
  t_c(f) \;=\; t_\mathrm{buf} + T\,
    \frac{\ln\!\left(f/f_1\right)}{\ln\!\left(f_2/f_1\right)},
  \label{eq:sweeptime}
\end{equation}
which is used to associate a time window with each excitation frequency.

Preprocessing consists of sorting the scan points along the cable axis,
removing dead channels (all-zero velocity), selecting the active shaker
channel by maximum energy, and correcting the known offset between the LDV
and shaker coordinate systems.

\paragraph{Band-pass filtering.}
% ldv_analysis/signal.py :: bandpass
Where a single excitation frequency $f_0$ is analysed, the velocities are
band-pass filtered with a zero-phase Butterworth filter of order 4
(applied forward and backward, \texttt{sosfiltfilt}, so the effective order
is 8 and the phase response is identically zero). The passband is
$f_0 \pm \SI{5}{Hz}$ (absolute bandwidth mode). Zero-phase filtering is
essential here: the coupling efficiency is a ratio of two simultaneously
filtered signals, and any filter-induced phase shift would bias it.

\paragraph{Integration to displacement.}
% ldv_analysis/signal.py :: integrate_fft
The LDV measures velocity; all coupling quantities are defined on
displacements. Integration is performed in the frequency domain. With
$V(f) = \mathcal{F}\{v\}(f)$ the discrete Fourier transform of the velocity,
the displacement spectrum follows from the differentiation property of the
Fourier transform,
\begin{equation}
  u(t) \;=\; \mathcal{F}^{-1}\!\left\{ \frac{V(f)}{i\,2\pi f} \right\}(t),
  \qquad f \ge f_\mathrm{min},
  \label{eq:integration}
\end{equation}
where the division is protected below $f_\mathrm{min} = \SI{0.5}{Hz}$
(frequencies below the floor are divided by $i\,2\pi f_\mathrm{min}$ instead)
to prevent the $1/f$ amplification of DC offset and low-frequency drift that
makes naive time-domain integration unstable. This is a linear, phase-exact
operation: a velocity component $v = \hat v\cos(2\pi f t)$ maps to
$u = \tfrac{\hat v}{2\pi f}\sin(2\pi f t)$, as expected.

\subsection{Geometry, Sag Determination and the Coupling Parameter $\Theta$}
\label{sec:proc:geometry}

\paragraph{Chord and axial projection.}
% ldv_analysis/geometry.py :: compute_chord, project_onto_chord
The cable spans a gap between two mounting points; the first and last scan
point inside the gap, $\mathbf{P}_1$ and $\mathbf{P}_2$, define the
\emph{chord}
\begin{equation}
  \hat{\mathbf{e}} \;=\; \frac{\mathbf{P}_2-\mathbf{P}_1}
                              {\lVert\mathbf{P}_2-\mathbf{P}_1\rVert},
  \qquad
  L \;=\; \lVert\mathbf{P}_2-\mathbf{P}_1\rVert .
  \label{eq:chord}
\end{equation}
Axial (in-line) displacement components are obtained by projection onto the
chord,
\begin{equation}
  u_\parallel(x_i, t) \;=\; \mathbf{u}(x_i,t)\cdot\hat{\mathbf{e}},
  \label{eq:projection}
\end{equation}
which is the displacement component a DAS fibre would sense. Projection onto
the chord (rather than taking the $x$-component) avoids mixing transverse
motion into the axial estimate when the cable is not perfectly aligned with
the coordinate axes.

\paragraph{Sag measurement.}
% ldv_analysis/geometry.py :: measure_sag (use_3d=True, use_parabola=True)
For every scan point in the gap, the perpendicular distance from the chord
line is computed in full 3-D:
\begin{equation}
  s_i = (\mathbf{P}_i - \mathbf{P}_1)\cdot\hat{\mathbf{e}},
  \qquad
  d_i = \bigl\lVert (\mathbf{P}_i-\mathbf{P}_1) - s_i\,\hat{\mathbf{e}}
        \bigr\rVert ,
  \label{eq:perpdist}
\end{equation}
where $s_i$ is the arc coordinate along the chord and $d_i$ the perpendicular
offset. Since the endpoints lie on the chord by construction
($d(0) = d(L) = 0$), a parabola \emph{constrained through both endpoints} is
fitted,
\begin{equation}
  d(s) \;\approx\; -a\, s\,(s - L),
  \qquad
  a \;=\; \frac{\sum_i \phi_i\, d_i}{\sum_i \phi_i^2},
  \quad \phi_i = s_i (s_i - L),
  \label{eq:parabolafit}
\end{equation}
a one-parameter linear least-squares problem. The midpoint sag is the vertex
value of the fitted parabola,
\begin{equation}
  w_0 \;=\; d(L/2) \;=\; \frac{a L^2}{4} \Big|_{a<0} ,
  \label{eq:sagvalue}
\end{equation}
and the root-mean-square residual of the fit,
$\sigma_{w_0} = \sqrt{\tfrac{1}{N}\sum_i (d_i - d(s_i))^2}$, is retained as
the sag measurement uncertainty (it captures outlier scan points and
deviations of the true shape from the parabolic approximation). The parabola
is the small-sag shape of a beam/cable under uniform load, so this fit is
consistent with the analytical model of
Section~\ref{sec:proc:geometry:theta}. If the fit yields $a \ge 0$ (no
physical downward arch), the maximum raw perpendicular distance is used
instead.

\paragraph{The coupling parameter $\Theta$.}
\label{sec:proc:geometry:theta}
% ldv_analysis/config.py :: theta_from_sag, theta_from_material
The bending-stress-relief model \citep{ProbstSimone} predicts that the
static strain-transfer efficiency of a suspended segment is
\begin{equation}
  \eta_\mathrm{stat} \;=\; \frac{\delta x_\ell}{\delta L}
                    \;=\; \frac{1}{1+\Theta},
  \qquad
  \Theta \;=\; \frac{EA}{L}\,\frac{(C_1 x_{f0})^2}{E I C_2},
  \label{eq:etatheta}
\end{equation}
where $x_{f0}$ is the gravity-induced midpoint sag,
$C_1 = \int_0^L (\phi')^2\,\mathrm{d}x$ and
$C_2 = \int_0^L (\phi'')^2\,\mathrm{d}x$ are the mode-shape constants of the
first symmetric clamped--clamped deflection shape
$\phi(x) = \tfrac{1}{2}\bigl[1-\cos(2\pi x/L)\bigr]$. Substituting the
tension-free equilibrium sag
\begin{equation}
  x_{f0} \;=\; w_0 \;=\; \frac{\rho A g L^4}{4\pi^4 E I}
  \label{eq:sagtheory}
\end{equation}
into Eq.~\eqref{eq:etatheta} gives the \emph{material form}
\begin{equation}
  \Theta_\mathrm{mat} \;=\; \frac{\rho^2 g^2 A^3 L^8}{128\,\pi^8 E^2 I^3},
  \label{eq:thetamaterial}
\end{equation}
while keeping the \emph{measured} sag $w_0$ gives the \emph{sag form}
\begin{equation}
  \Theta_\mathrm{sag} \;=\; \frac{A\, w_0^2}{8 I}
  \;\overset{\text{circular}}{=}\; \frac{1}{2}\left(\frac{w_0}{r}\right)^2 ,
  \label{eq:thetasag}
\end{equation}
where the last equality holds for a circular cross-section with
$A = \pi r^2$ and $I = \pi r^4/4$. For the rectangular ribbon cable
(width $w$, height $h$, bending about the weak axis) $A = wh$ and
$I = w h^3/12$ are used, so Eq.~\eqref{eq:thetasag} correctly generalises
beyond circular cables. The sag form is preferred throughout this work
because it absorbs the actual (pretension-affected) deployment geometry; the
material form is reported for comparison. Young's modulus enters as the mean
of three independent tensile tests per cable; $\Theta_\mathrm{mat}$ is
evaluated once per test value and reported as mean~$\pm$~standard deviation.

\subsection{Strain and Elongation Estimation}
\label{sec:proc:strain}

The coupling efficiency compares the \emph{axial elongation of the cable}
with the \emph{elongation imposed at its boundaries}. Three independent
estimators of the cable elongation are implemented --- M1 per-segment 3-D
arc-length, M2 spatial gradient and M3 Fourier-domain gradient --- and their
agreement is a built-in consistency check of the processing.

\subsubsection{Method 1: Per-Segment 3-D Arc-Length Elongation}
\label{sec:proc:strain:m1}
% ldv_analysis/strain.py :: strain_3d_arclength
Method 1 estimates the elongation of each segment between two adjacent scan
points directly from their 3-D displacement vectors, without projecting the
whole field onto the global chord first. Let
$\mathbf{r}_i = \mathbf{P}_{i+1} - \mathbf{P}_i$ be the static segment
vector, $L_{0,i} = \lVert\mathbf{r}_i\rVert$ its rest length and
$\hat{\mathbf{e}}_i = \mathbf{r}_i / L_{0,i}$ its \emph{local} unit vector.
The deformed segment length is
\begin{equation}
  L_i(t) \;=\; \bigl\lVert \mathbf{r}_i
      + \Delta\mathbf{u}_i(t) \bigr\rVert ,
  \qquad
  \Delta\mathbf{u}_i(t) = \mathbf{u}_{i+1}(t) - \mathbf{u}_i(t).
\end{equation}
Because the dynamic displacements ($\sim\si{\micro\meter}$) are many orders
of magnitude smaller than the segment lengths ($\sim\si{\milli\meter}$), the
square root is linearised by a first-order Taylor expansion:
\begin{equation}
  L_i(t) - L_{0,i}
  \;=\; L_{0,i}\sqrt{1
        + \frac{2\,\mathbf{r}_i\!\cdot\!\Delta\mathbf{u}_i}{L_{0,i}^2}
        + \frac{\lVert\Delta\mathbf{u}_i\rVert^2}{L_{0,i}^2}} - L_{0,i}
  \;\approx\;
  \hat{\mathbf{e}}_i \cdot \Delta\mathbf{u}_i(t)
  \;\equiv\; \delta d_i(t).
  \label{eq:strainM1}
\end{equation}
Geometrically: \emph{the relative displacement of the two points is projected
onto the local segment direction} $\hat{\mathbf{e}}_i$ — only the component
of relative motion along the segment changes its length to first order; the
perpendicular component merely rotates it (a second-order length effect that
the linearisation deliberately drops, consistent with the small-strain
regime). For a sagged cable the $\hat{\mathbf{e}}_i$ follow the local cable
direction, so Method 1 measures the \emph{arc-length} change, which is
exactly the quantity a strained optical fibre integrates. The local strain
estimate of segment $i$ is $\varepsilon_i = \delta d_i / L_{0,i}$, and the
total cable elongation is the telescoping sum over the gap,
\begin{equation}
  \delta x_\ell(t) \;=\; \sum_{i\,\in\,\mathrm{gap}} \delta d_i(t).
  \label{eq:deltaxl}
\end{equation}
Method 1 is the primary estimator in this work: it makes no assumption
about the smoothness of the displacement field between scan points and is
insensitive to the chord-projection error for sagged geometries.

\subsubsection{Method 2: Spatial Gradient of the Axial Displacement}
% ldv_analysis/strain.py :: strain_spatial_gradient
Axial strain is the spatial derivative of the axial displacement along the
cable,
\begin{equation}
  \varepsilon(s,t) \;=\; \frac{\partial\, u_\parallel(s,t)}{\partial s},
  \label{eq:strainM2def}
\end{equation}
evaluated with second-order central finite differences on the non-uniform
scan-point grid. For interior point $i$ with spacings
$\Delta s_- = s_i - s_{i-1}$ and $\Delta s_+ = s_{i+1} - s_i$:
\begin{equation}
  \varepsilon_i \;=\;
    u_{\parallel,i+1}\,\frac{\Delta s_-}{\Delta s_+(\Delta s_+ + \Delta s_-)}
  + u_{\parallel,i}\,\frac{\Delta s_+ - \Delta s_-}{\Delta s_+ \Delta s_-}
  - u_{\parallel,i-1}\,\frac{\Delta s_+}{\Delta s_-(\Delta s_+ + \Delta s_-)} .
  \label{eq:strainM2}
\end{equation}
This reduces to the familiar $(u_{i+1}-u_{i-1})/(2\Delta s)$ for uniform
spacing and is exact for locally quadratic displacement fields. The total
elongation follows by trapezoidal integration over the gap,
$\delta x_\ell^{(2)}(t) = \int_\mathrm{gap} \varepsilon(s,t)\,\mathrm{d}s$.

\subsubsection{Method 3: Fourier-Domain Spatial Gradient}
% ldv_analysis/strain.py :: strain_fourier_gradient
As a cross-check with different numerical error characteristics, the axial
displacement is resampled onto a uniform grid $s_k$ (cubic interpolation)
and differentiated spectrally using the Fourier differentiation theorem,
\begin{equation}
  \varepsilon(s,t) \;=\;
  \mathcal{F}_s^{-1}\bigl\{\, i k\, \mathcal{F}_s\{u_\parallel\}\,\bigr\},
  \label{eq:strainM3}
\end{equation}
with $k$ the spatial angular wavenumber. Spectral differentiation is exact
for band-limited periodic signals; since the cable displacement is not
periodic over the gap, Gibbs-type artefacts can appear near the endpoints,
which is why only gap-interior stations enter the elongation integral. In
practice Method 3 agrees with Methods 1--2 in the smooth interior and serves
to expose discretisation artefacts of the finite-difference estimate.

\subsubsection{Reference Elongations}
% ldv_analysis/freqdomain.py :: build_coupling_signals (x_source)
Three references for the imposed boundary elongation $\delta L(t)$ are
available:
\begin{itemize}
  \item \textbf{Endpoint pair} (default): the span change measured on the
    cable itself,
    $\delta L = u_\parallel(\mathbf{P}_2) - u_\parallel(\mathbf{P}_1)$.
    It shares its endpoints with $\delta x_\ell$, so the efficiency isolates
    bending stress relief and is immune to cross-instrument calibration.
  \item \textbf{Shaker}: the chord-projected shaker-head displacement
    $\delta L = \mathbf{u}_\mathrm{sh}\cdot\hat{\mathbf{e}}$ (sign-corrected
    so that positive means stretching), i.e.\ the imposed ground motion.
  \item \textbf{Shaker-to-end}: the elongation of the sub-span between the
    shaker and the adjacent endpoint sensor, used to diagnose losses at the
    first contact.
\end{itemize}

\subsection{Time-Domain Coupling Efficiency (Quality Control)}
\label{sec:proc:timedomain}
% ldv_analysis/analysis.py :: per_frequency_qc

For inspection at a single excitation frequency $f_0$, the record is windowed
to $t_c(f_0) \pm n_\mathrm{cyc}/f_0$ (Eq.~\ref{eq:sweeptime},
$n_\mathrm{cyc}=5$), band-passed (Section~\ref{sec:proc:preprocessing}),
integrated to displacement, and the instantaneous efficiency is formed:
\begin{equation}
  \eta(t) \;=\; \frac{\delta x_\ell(t)}{\delta L(t)},
  \qquad
  \eta_\mathrm{med} = \operatorname{median}_t\, \eta(t)
  \quad\text{over } |\delta L(t)| > 0.05\,\max|\delta L| .
  \label{eq:etatime}
\end{equation}
The amplitude gate excludes the zero crossings of the reference, where the
ratio is ill-conditioned; the median suppresses the remaining outliers. A
phase-insensitive variant replaces both signals by their analytic-signal
envelopes (Section~\ref{sec:proc:amplitude}),
$\eta_\mathrm{env}(t) = |\mathcal{H}\{\delta x_\ell\}| /
|\mathcal{H}\{\delta L\}|$: comparing $\eta_\mathrm{med}$ and
$\eta_\mathrm{env}$ distinguishes whether a low efficiency is caused by an
amplitude deficit or by a phase lag between cable and reference.

These time-domain values are used for quality control and for illustrating
the mechanism at single frequencies; all summary efficiencies in the results
chapter are computed in the frequency domain
(Section~\ref{sec:proc:freqdomain}), which uses the whole record at once and
is considerably more stable.

\subsection{Frequency-Domain Transfer Functions}
\label{sec:proc:freqdomain}
% ldv_analysis/freqdomain.py

\subsubsection{Signals}
From the full (unfiltered) sweep record, integrated to displacement, three
signals are built:
\begin{align}
  X(t) &= \delta L(t) && \text{input: boundary span change,} \notag\\
  Y_1(t) &= \delta x_\ell(t) && \text{output 1: cable elongation
                                (Eq.~\ref{eq:deltaxl}),} \notag\\
  Y_2(t) &= u_{\perp,\mathrm{mid}}(t) && \text{output 2: midpoint deflection
                                along the static sag direction.} \notag
\end{align}
$H_\eta(f) = Y_1/X$ is the dynamic strain-transfer efficiency;
$H_\mathrm{mid}(f) = Y_2/X$ measures the intensity of the bending mode, whose
resonance structure explains the frequency dependence of $H_\eta$
\citep{ProbstSimone}.

\subsubsection{Welch Spectral Estimation}
% ldv_analysis/freqdomain.py :: welch_spectra
The record is divided into $K$ segments of length
$n_\mathrm{seg} = 2 f_s = 10\,000$ samples (\SI{2}{s}, frequency resolution
$\Delta f = \SI{0.5}{Hz}$) with 50\,\% overlap, each tapered with a Hann
window $w[n]$. With $X_k(f)$, $Y_k(f)$ the windowed DFTs of segment $k$,
the auto- and cross-spectral densities are the segment averages
\begin{equation}
  S_{XX}(f) = \frac{1}{K}\sum_{k=1}^{K} |X_k(f)|^2 , \qquad
  S_{YY}(f) = \frac{1}{K}\sum_{k=1}^{K} |Y_k(f)|^2 , \qquad
  S_{XY}(f) = \frac{1}{K}\sum_{k=1}^{K} X_k^*(f)\, Y_k(f),
  \label{eq:welch}
\end{equation}
(up to the standard window normalisation, which cancels in every ratio used
below). The averaging over $K \approx 11$ segments reduces the variance of
each spectral estimate by $\approx 1/K$ — this is what makes the Welch-based
transfer functions far smoother than any single-FFT estimate. The
\emph{ordinary coherence}
\begin{equation}
  \gamma^2(f) \;=\; \frac{|S_{XY}(f)|^2}{S_{XX}(f)\, S_{YY}(f)}
  \;\in\; [0, 1]
  \label{eq:coherence}
\end{equation}
measures the fraction of output power linearly explained by the input;
$\gamma^2 < 1$ indicates noise, nonlinearity, or unmeasured inputs.

\subsubsection{FRF Estimators: Direct, $H_1$, $H_2$, $H_v$}
\label{sec:proc:estimators}
% ldv_analysis/freqdomain.py :: direct_transfer, estimate_H

\paragraph{Direct estimate.}
The conceptually simplest frequency response estimate is the ratio of the
full-record windowed Fourier transforms,
\begin{equation}
  H_\mathrm{direct}(f) \;=\;
  \frac{\mathcal{F}\{w\,Y\}(f)}{\mathcal{F}\{w\,X\}(f)} ,
  \label{eq:Hdirect}
\end{equation}
evaluated on the dense grid $\Delta f_\mathrm{d} = f_s/N = \SI{0.083}{Hz}$.
Every frequency bin is a single realisation: the estimate is unbiased in the
noise-free limit but has no variance reduction, so amplitude and phase
scatter strongly wherever excitation or response is weak.

Comparing Eq.~\eqref{eq:Hdirect} directly against $H_1$ would conflate two
distinct effects — the estimator algebra and the amount of spectral averaging.
To separate them, the direct spectra are smoothed along the \emph{frequency}
axis (Daniell smoothing) rather than over time segments: with a normalised
kernel $k$ of $m$ bins,
\begin{equation}
  \tilde S_{AB}(f_i) = \sum_{j} k_j\, A^*(f_{i+j})\, B(f_{i+j}),
  \qquad
  H_\mathrm{direct}^{\mathrm{sm}}(f) = \frac{\tilde S_{XY}(f)}{\tilde S_{XX}(f)} .
  \label{eq:Hdirectsmooth}
\end{equation}
Smoothing the products \emph{before} the division makes the estimate
input-power weighted, so bins with weak $|X|$ contribute little instead of
diverging; this also removes the positive bias that the unweighted ratio
carries at low SNR. A kernel with $\sum_j k_j = 1$ averages
$n_\mathrm{eff} = 1/\sum_j k_j^2$ independent bins, which sets both the
variance reduction and the effective resolution bandwidth. The kernel width is
therefore chosen so that $n_\mathrm{eff}\,\Delta f_\mathrm{d}$ equals the
\emph{equivalent noise bandwidth} of the Welch estimate,
\begin{equation}
  B_\mathrm{ENBW} = f_s \frac{\sum_n w[n]^2}{\left(\sum_n w[n]\right)^2}
  = 1.5\,\Delta f = \SI{0.75}{Hz}
  \quad\text{(Hann)},
  \label{eq:enbw}
\end{equation}
i.e.\ $n_\mathrm{eff} \approx 9$ dense bins. Note that
$B_\mathrm{ENBW} > \Delta f$: the Hann taper spreads each spectral line over
1.5 bins, so a Welch estimate does not resolve to its nominal bin spacing, and
matching $B_\mathrm{ENBW}$ rather than $\Delta f$ is what puts the two
estimates on an equal footing. With this choice the smoothed direct estimate
and $H_1$ agree to within a few percent in $|H|$ while retaining the dense
frequency grid; the residual difference is attributable to the estimator
algebra alone. Both the raw and the smoothed direct curves are shown, so the
figure separates the effect of averaging (raw $\to$ smoothed) from the effect
of the estimator (smoothed $\to$ $H_1$).

\paragraph{$H_1$ — noise on the output.}
Assume the measured output contains additive noise uncorrelated with the
input, $Y = H X + N$. Minimising the output residual over the segment
ensemble,
$\min_H \sum_k |Y_k - H X_k|^2$, gives the least-squares solution
\begin{equation}
  H_1(f) \;=\; \frac{S_{XY}(f)}{S_{XX}(f)} .
  \label{eq:H1}
\end{equation}
Since $\mathbb{E}\{X^* N\} = 0$, the cross-spectrum averages the output noise
away and $H_1$ is unbiased under output noise. If the \emph{input} channel is
noisy, $S_{XX}$ is inflated and $|H_1|$ \emph{under}estimates the true FRF.

\paragraph{$H_2$ — noise on the input.}
Assume instead noise on the input, $X = X_0 + M$ with $Y = H X_0$.
Minimising the input-referenced residual
$\min_H \sum_k |X_k - Y_k / H|^2$ yields
\begin{equation}
  H_2(f) \;=\; \frac{S_{YY}(f)}{S_{YX}(f)}
        \;=\; \frac{S_{YY}(f)}{S_{XY}^*(f)} .
  \label{eq:H2}
\end{equation}
$H_2$ is unbiased under input noise but \emph{over}estimates $|H|$ when the
output is noisy. The two estimators bracket the true FRF,
\begin{equation}
  |H_1(f)| \;\le\; |H(f)| \;\le\; |H_2(f)|,
  \qquad
  \frac{|H_1(f)|}{|H_2(f)|} = \gamma^2(f),
  \label{eq:bracket}
\end{equation}
so their separation is itself a coherence display: the estimators agree
exactly where $\gamma^2 = 1$ and diverge where the data are noisy.

\paragraph{$H_v$ — noise on both channels (total least squares).}
When both channels carry comparable noise — realistic here, since input and
output are both derived from LDV velocities through the same processing —
the total-least-squares estimator distributes the residual over both
channels. Writing the spectral sample covariance of the vector $(X, Y)$ per
frequency bin,
\begin{equation}
  \mathbf{C}(f) \;=\;
  \begin{pmatrix} S_{XX} & S_{XY}^* \\ S_{XY} & S_{YY} \end{pmatrix},
\end{equation}
the TLS solution is the direction of the eigenvector belonging to the
\emph{smallest} eigenvalue $\lambda_-$ of $\mathbf{C}$ (the noise
eigenvalue). For the $2\times 2$ Hermitian matrix this eigenvalue is
\begin{equation}
  \lambda_-(f) \;=\; \frac{S_{XX}+S_{YY}}{2}
    - \sqrt{\left(\frac{S_{XX}-S_{YY}}{2}\right)^{\!2} + |S_{XY}|^2},
\end{equation}
and solving the eigenvector relation for the FRF gives the closed form
\begin{equation}
  H_v(f) \;=\; \frac{S_{YY}(f) - \lambda_-(f)}{S_{XY}^*(f)} .
  \label{eq:Hv}
\end{equation}
One verifies directly that in the noise-free case ($\lambda_- = 0$,
$S_{YY} = |H|^2 S_{XX}$, $S_{XY} = H S_{XX}$) Eq.~\eqref{eq:Hv} returns
exactly $H$, and that $|H_1| \le |H_v| \le |H_2|$ holds bin-by-bin. All
three Welch estimators share the same phase,
$\arg H_1 = \arg H_2 = \arg H_v = \arg S_{XY}$, so the estimator choice
affects only the magnitude.

\paragraph{Phase presentation.}
Phases are reported \emph{wrapped} to $(-180^\circ, 180^\circ]$. Unwrapped
phase curves accumulate spurious $\pm 360^\circ$ jumps whenever
low-coherence bins randomise the phase, which makes them diverge without
physical meaning; the wrapped display is immune to this and the physically
relevant transition (the $\approx 180^\circ$ flip across the bending
resonance) remains fully visible.

\subsubsection{Band-Averaged Efficiency and Coherence Gating}
% ldv_analysis/freqdomain.py :: band_mean_eta
The quasi-static efficiency reported per configuration is the coherence-gated
band average
\begin{equation}
  \bar\eta \;=\; \bigl\langle\, |H_\eta(f)| \,\bigr\rangle_{f \in B},
  \qquad
  B = \{ f \in [f_\mathrm{min}, f_\mathrm{max}] :
         \gamma^2_\eta(f) \ge 0.7 \},
  \label{eq:bandeta}
\end{equation}
with $[f_\mathrm{min}, f_\mathrm{max}] = [5, 25]\,\si{Hz}$ well below the
first resonance. The gate removes noise-dominated bins that would otherwise
bias the average; the standard deviation over $B$ is reported as the error
bar.

\subsubsection{Single-Frequency Cross-Check and Resonance Characterisation}
% ldv_analysis/freqdomain.py :: stft_point_transfer, fit_lorentzian
As an independent validation, point estimates
$H(f_0) = \mathcal{F}\{w\,Y\}(f_0)\,/\,\mathcal{F}\{w\,X\}(f_0)$ are computed
from short Hann-tapered windows centred on the sweep arrival $t_c(f_0)$ and
overlaid on the Welch curves.

The fundamental resonance $f_1$ is identified as the global maximum of
$|H_\mathrm{mid}(f)|$; a driven single-degree-of-freedom magnitude response
\begin{equation}
  |H(f)| \;=\; \frac{G\, f_1^2}
   {\sqrt{\left(f_1^2 - f^2\right)^2 + \left(f_1 f / Q\right)^2}}
  \label{eq:sdof}
\end{equation}
is fitted in a window around the peak, yielding the quality factor $Q$ and
damping ratio $\zeta = 1/(2Q)$. The theoretical clamped--clamped fundamental
\begin{equation}
  f_1^\mathrm{th} \;=\; \frac{\alpha_1^2}{2\pi L^2}\sqrt{\frac{EI}{\rho A}},
  \qquad \alpha_1^2 = 4.730^2 \approx 22.37,
  \label{eq:f1theory}
\end{equation}
is quoted for comparison (cf.~Section~\ref{sec:proc:modal}).

\subsection{Amplitude Extraction}
\label{sec:proc:amplitude}
% ldv_analysis/amplitude.py

Beyond ratios, the absolute displacement and strain amplitudes are needed as
reference values for the numerical simulations. Two extraction routes are
implemented and cross-validated.

\paragraph{Time domain: analytic-signal envelope.}
For a windowed, band-passed trace $x(t)$ the analytic signal is
$x_a(t) = x(t) + i\,\mathcal{H}\{x\}(t)$ with $\mathcal{H}$ the Hilbert
transform; its magnitude $|x_a(t)|$ is the instantaneous envelope. The
representative amplitude is the envelope median (robust against the
filter-edge droop at the window boundaries); envelope maximum and
$\sqrt{2}\,\mathrm{RMS}$ variants are available for comparison.

\paragraph{Frequency domain: band-energy read-out.}
The velocities are integrated once per dataset; at each target frequency
$f_0$ the displacement signal is windowed around $t_c(f_0)$, Hann-tapered
and Fourier transformed. Because the log-sweep smears the tone over several
bins (especially at low frequency, where the window spans many seconds), the
amplitude is recovered from the \emph{energy in the analysis band}
$B = [f_0 - 5, f_0 + 5]\,\si{Hz}$ rather than from a single bin. For a tone
of peak amplitude $A$ under a window $w[n]$, Parseval's theorem gives
\begin{equation}
  A \;=\; 2\sqrt{\frac{\sum_{k \in B} |X_k|^2}{N \sum_n w[n]^2}} ,
  \label{eq:bandenergy}
\end{equation}
which is exact for an isolated in-band tone. This route avoids the
per-frequency band-pass filtering of the full record and is roughly two
orders of magnitude faster, enabling dense frequency grids; it agrees with
the envelope route to within a few percent below $\approx\SI{100}{Hz}$.

\subsection{Extension--Compression Asymmetry}
\label{sec:proc:asymmetry}
% ldv_analysis/asymmetry.py

A sagged segment stiffens when its endpoints are pulled apart (the sag is
straightened) and softens when they are pushed together (the sag deepens);
numerical simulations show correspondingly different responses for the two
deformation directions. A perfectly linear system responds
mirror-symmetrically to the two half-cycles of a sinusoidal drive, so any
measurable asymmetry is a signature of this geometric nonlinearity. Three
complementary detectors are evaluated per excitation frequency:

\begin{enumerate}
  \item \textbf{Half-cycle conditional efficiency.} The samples of the
  reference elongation are split by sign — extension
  ($\delta L > 0$, endpoints separating) versus compression
  ($\delta L < 0$) — with a $25\,\%$-of-maximum amplitude gate excluding the
  zero crossings. The efficiencies
  \begin{equation}
    \eta_\mathrm{ext} = \operatorname*{median}_{\delta L(t) > \delta_g}
      \frac{\delta x_\ell(t)}{\delta L(t)},
    \qquad
    \eta_\mathrm{comp} = \operatorname*{median}_{\delta L(t) < -\delta_g}
      \frac{\delta x_\ell(t)}{\delta L(t)}
    \label{eq:etahalfcycle}
  \end{equation}
  coincide for a linear system; $\eta_\mathrm{ext} \ne \eta_\mathrm{comp}$
  quantifies the directional coupling difference.

  \item \textbf{Peak asymmetry.} From the mean positive-peak amplitude
  $A_+$ and mean negative-peak magnitude $A_-$ of the cable elongation
  $\delta x_\ell(t)$:
  \begin{equation}
    \alpha \;=\; \frac{A_+ - A_-}{A_+ + A_-} \;\in\; [-1, 1],
    \label{eq:peakasym}
  \end{equation}
  with $\alpha > 0$ meaning the cable elongates further during extension
  half-cycles than it shortens during compression.

  \item \textbf{Even-harmonic distortion.} Expanding the response in powers
  of the drive, $y = c_1 x + c_2 x^2 + c_3 x^3 + \dots$ with
  $x = A\cos\omega t$, the quadratic term produces
  $c_2 A^2 \cos^2\omega t = \tfrac{c_2 A^2}{2}(1 + \cos 2\omega t)$ — energy
  at \emph{twice} the drive frequency — whereas the cubic (symmetric) term
  appears at $3\omega$. Since a sign-asymmetric stiffness is exactly a
  quadratic term, the ratio
  \begin{equation}
    R_2 \;=\; \frac{|Y(2 f_0)|}{|Y(f_0)|}
    \label{eq:h2ratio}
  \end{equation}
  isolates the asymmetric part of the response independently of the
  reference channel; $R_3 = |Y(3f_0)|/|Y(f_0)|$ is reported alongside to
  attribute nonlinearity to the symmetric (softening/stiffening) branch.
\end{enumerate}

For the linearity data set (Section~\ref{sec:proc:linearity}) the half-cycle
efficiencies are additionally evaluated in bins of the instantaneous drive
amplitude, revealing whether an asymmetry develops progressively with
amplitude (e.g.\ slack forming on the compression half-cycle only).

\subsection{Modal Analysis}
\label{sec:proc:modal}
% ldv_analysis/modal.py

\subsubsection{Detection Fingerprint and Spectral Smoothing}
% modal.spectral_fingerprint
Transverse resonances are detected on a \emph{fingerprint}: per-sensor
spectra of the transverse velocities $v_y, v_z$ over the gap, collapsed
across sensors into the total-power detection signal
\begin{equation}
  P(f) \;=\; \sum_{s\,\in\,\mathrm{gap}}
      \Bigl( |V_y(f; s)|^2 + |V_z(f; s)|^2 \Bigr).
  \label{eq:fingerprint}
\end{equation}
The per-sensor spectra are Welch estimates
($n_\mathrm{seg} = 2 f_s$, $\Delta f = \SI{0.5}{Hz}$, Hann, 50\,\% overlap,
$K \approx 11$ averages) rather than raw single-record FFTs: the
$\mathrm{var} \propto 1/K$ averaging suppresses the bin-to-bin fluctuations
that previously caused the peak picker to fire on noise wiggles. An optional
moving-average (boxcar) smoothing over a configurable bandwidth can be
applied on top. Resonances are then the peaks of $P(f)$ (normalised to its
in-band maximum) exceeding a prominence threshold with a minimum mutual
spacing (\texttt{scipy.signal.find\_peaks}).

\paragraph{Input-bias correction.}
% modal.shaker_input_spectrum, normalize_by_input
The shaker head does not excite purely axially: its own recording shows
frequency-dependent energy in the transverse components. A spectral peak in
the cable response may therefore reflect a peak in the \emph{excitation}
rather than a cable resonance. Two diagnostics address this. First, the
input directionality is quantified by the per-frequency power fractions of
the shaker-head components,
\begin{equation}
  \phi_c(f) \;=\; \frac{|V_c^\mathrm{sh}(f)|^2}
   {\sum_{c' \in \{x,y,z\}} |V_{c'}^\mathrm{sh}(f)|^2},
  \qquad c \in \{x, y, z\}.
  \label{eq:directionality}
\end{equation}
Second, the detection signal can be normalised by the total shaker input
power, $P_\mathrm{norm}(f) = P(f) / P_\mathrm{sh}(f)$ — an FRF-like quantity
in which input coloration cancels: peaks that survive the normalisation are
genuine output-per-input amplifications (resonances), while peaks that
disappear were inherited from the excitation spectrum.

\subsubsection{Complex Mode-Shape Extraction}
% modal.extract_mode_shape
For each detected resonance $f_n$, the transverse velocities are band-passed
to $f_n(1 \pm 0.1)$, integrated to displacement, and the complex Fourier
coefficient at the bin nearest $f_n$ is read per sensor:
\begin{equation}
  \varphi_y(s) = U_y(f_n; s), \qquad \varphi_z(s) = U_z(f_n; s).
  \label{eq:modeshape}
\end{equation}
The complex-valued shapes retain the relative phase between sensors, so
travelling-wave content and sign flips across nodes are preserved.

\subsubsection{Euler--Bernoulli Reference Modes}
% modal.clamped_clamped_mode
The gap is modelled as a doubly-clamped uniform beam. Free vibration of an
Euler--Bernoulli beam obeys $EI\, w'''' = \rho A\, \omega^2 w$; with clamped
boundary conditions $w(0)=w'(0)=w(L)=w'(L)=0$, non-trivial solutions exist
only at the roots $\alpha_n$ of the frequency equation
\begin{equation}
  \cos\alpha \cosh\alpha \;=\; 1,
  \qquad
  \alpha_n = 4.730,\ 7.853,\ 10.996,\ 14.137,\ 17.279,\ \dots
  \label{eq:freqequation}
\end{equation}
with natural frequencies and mode shapes
\begin{equation}
  f_n \;=\; \frac{\alpha_n^2}{2\pi L^2} \sqrt{\frac{EI}{\rho A}},
  \qquad
  \varphi_n(\xi) = \bigl[\cosh\alpha_n\xi - \cos\alpha_n\xi\bigr]
   - \sigma_n \bigl[\sinh\alpha_n\xi - \sin\alpha_n\xi\bigr],
  \label{eq:ebmodes}
\end{equation}
where $\xi = x/L \in [0,1]$ and
$\sigma_n = (\cosh\alpha_n - \cos\alpha_n)/(\sinh\alpha_n - \sin\alpha_n)$.
The frequency ratios $f_n / f_1 = (\alpha_n/\alpha_1)^2 =
1,\ 2.76,\ 5.40,\ 8.93,\ \dots$ provide a material-free consistency check of
the mode ordering.

\subsubsection{Modal Assurance Criterion}
% modal.mac
Each measured shape is assigned to a theoretical mode by the Modal Assurance
Criterion,
\begin{equation}
  \operatorname{MAC}(\boldsymbol{\varphi}_m, \boldsymbol{\varphi}_t)
  \;=\;
  \frac{\bigl|\boldsymbol{\varphi}_m^{\mathsf{H}}
              \boldsymbol{\varphi}_t\bigr|^2}
       {\bigl(\boldsymbol{\varphi}_m^{\mathsf{H}}\boldsymbol{\varphi}_m\bigr)
        \bigl(\boldsymbol{\varphi}_t^{\mathsf{H}}\boldsymbol{\varphi}_t\bigr)}
  \;\in\; [0, 1],
  \label{eq:mac}
\end{equation}
where $\boldsymbol{\varphi}_m$ is the measured complex shape (the dominant
polarization component, sampled at the sensor positions),
$\boldsymbol{\varphi}_t$ the theoretical shape $\varphi_n(\xi_i)$ evaluated
at the same positions, and $^{\mathsf{H}}$ the conjugate transpose. The MAC
is the squared modulus of the normalised inner product — a
\emph{shape-correlation} coefficient: it equals 1 when the vectors are
collinear (identical shape up to any complex scale), and 0 when they are
orthogonal. Because a global complex factor cancels, the measured shape
needs no phase alignment or amplitude scaling before comparison, and a
$180^\circ$ phase flip across a node is handled correctly. The measured
resonance is assigned the mode order with the largest MAC; low best-MAC
values flag shapes that no beam mode explains (e.g.\ rigid-body-like motion
from mount compliance or shaker cross-talk).

\subsubsection{Polarization Analysis}
% modal.polarization_ellipse
At each sensor the pair $(\varphi_y, \varphi_z)$ defines an ellipse traced
by $\Re\{(\varphi_y, \varphi_z)\, e^{i\omega t}\}$ in the $y$--$z$ plane.
Decomposing into counter-rotating circular components
\begin{equation}
  A_+ = \tfrac{1}{2}\bigl(\varphi_y + i\,\varphi_z\bigr),
  \qquad
  A_- = \tfrac{1}{2}\bigl(\varphi_y^* + i\,\varphi_z^*\bigr)
  \quad\Longrightarrow\quad
  \text{major} = |A_+| + |A_-|,\quad
  \text{minor} = \bigl||A_+| - |A_-|\bigr|,
  \label{eq:ellipse}
\end{equation}
with axis ratio $\mathrm{minor}/\mathrm{major}$ (0 = linear, 1 = circular),
orientation $\theta = \tfrac{1}{2}(\arg A_+ + \arg A_-)$ and handedness
$\operatorname{sign}\Im(\varphi_y \varphi_z^*)$. A resonance is classified
linear / elliptical / circular by the amplitude-weighted majority over the
significant sensors.

\subsection{Linearity Analysis}
\label{sec:proc:linearity}
% ldv_analysis/analysis.py :: linearity_qc

A dedicated test signal probes amplitude dependence: eight discrete
frequencies (5--400\,\si{Hz}), each driven for \SI{3.5}{s} with a linear
amplitude ramp $0 \to 1$ followed by a \SI{0.5}{s} fade. Each segment is
extracted in full, band-passed ($f_0 \pm \SI{5}{Hz}$), integrated, and the
efficiency $\eta(t)$ (Eq.~\ref{eq:etatime}) is evaluated against the
\emph{instantaneous drive amplitude}, defined as the analytic-signal envelope
$a(t) = |\mathcal{H}\{\delta L\}(t)|$. Because the ramp sweeps amplitude
monotonically, plotting the binned median of $\eta$ against $a$ directly
tests linearity: a flat curve means amplitude-independent coupling, a slope
or hysteresis loop indicates nonlinearity (e.g.\ frictional slip at the
mounts above a threshold amplitude). A complementary slip indicator divides
the envelope of a single cable-sensor velocity by the envelope of the shaker
reference; for slip-free coupling this ratio stays constant along the ramp.
The frequency-domain cross-check estimates $H_1$
(Eq.~\ref{eq:H1}) on each mono-frequency segment and compares
$|H_1(f_0)|$ with the time-domain median efficiency.

\subsection{Uncertainty Treatment}
\label{sec:proc:uncertainty}

\begin{itemize}
  \item \textbf{Young's modulus:} three tensile tests per cable;
    $\Theta_\mathrm{mat}$ is evaluated per test and reported as mean $\pm$
    standard deviation. Since $\Theta_\mathrm{mat} \propto E^{-2}$, a
    relative modulus error $\epsilon_E$ propagates as
    $2\,\epsilon_E$ into $\Theta$.
  \item \textbf{Gap length:} the mounting-gap reading error is
    $\sigma_L = \SI{0.5}{cm}$. The material form is extremely sensitive,
    $\Theta_\mathrm{mat} \propto L^8 \Rightarrow
    \sigma_\Theta/\Theta = 8\,\sigma_L/L$ (40\,\% at $L=\SI{10}{cm}$!),
    which is the quantitative reason the sag-based
    $\Theta_\mathrm{sag}$ — independent of $L$ — is preferred.
  \item \textbf{Sag:} the RMS residual of the constrained parabola fit
    (Eq.~\ref{eq:parabolafit}) serves as $\sigma_{w_0}$; with
    $\Theta_\mathrm{sag} \propto w_0^2$,
    $\sigma_\Theta/\Theta = 2\,\sigma_{w_0}/w_0$.
  \item \textbf{Efficiency:} the standard deviation of $|H_\eta|$ over the
    coherence-gated quasi-static band (Eq.~\ref{eq:bandeta}); the
    $H_1$/$H_2$ bracket (Eq.~\ref{eq:bracket}) provides an additional
    systematic error envelope.
\end{itemize}

% ─── Bibliography hooks ─────────────────────────────────────────────────────
% \citep{ProbstSimone}  : S. Probst et al., "Unburied DAS on the Moon:
%                         Modeling Ground-to-Cable Strain Transfer" (draft).
% Probst et al. 2026    : Controlled Source DAS Coupling Tests, Earth and
%                         Space Science 13(3), e2025EA004817.
% Welch 1967            : P. Welch, IEEE Trans. Audio Electroacoust. 15(2).
% Allemang 2003         : R. Allemang, "The Modal Assurance Criterion --
%                         Twenty Years of Use and Abuse", Sound & Vibration.
% Bendat & Piersol 2010 : Random Data: Analysis and Measurement Procedures
%                         (H1/H2/Hv estimators, coherence).

% Required packages: amsmath, amssymb, siunitx, booktabs (for the tables).
% Every equation cites the implementing function in ldv_analysis/ in a margin
% comment so the text can be checked against the code.
% ═══════════════════════════════════════════════════════════════════════════

\section{Signal Processing and Data Analysis}
\label{sec:processing}

All quantitative results in this thesis are derived from the raw Laser
Doppler Vibrometer (LDV) recordings through the processing chain described in
this chapter. The chain is implemented in the Python package
\texttt{ldv\_analysis}; each subsection below corresponds to one module of
that package, and equations are stated exactly as implemented.

\subsection{Measurement Data and Preprocessing}
\label{sec:proc:preprocessing}

Each measurement consists of two synchronous recordings: the scanning LDV
measures the three-component velocity field
$\mathbf{v}(x_i, t) = (v_x, v_y, v_z)$ at $N_s$ scan points along the cable,
and a reference channel records the three-component velocity of the shaker
head, $\mathbf{v}_\mathrm{sh}(t)$. The excitation is a logarithmic sweep from
$f_1 = \SI{1}{Hz}$ to $f_2 = \SI{500}{Hz}$ over $T = \SI{10}{s}$ (after a
\SI{1}{s} buffer), sampled at $f_s = \SI{5}{kHz}$ with $N = 60\,000$ samples.
The instantaneous sweep frequency reaches $f$ at
% ldv_analysis/config.py :: sweep_time_of_frequency
\begin{equation}
  t_c(f) \;=\; t_\mathrm{buf} + T\,
    \frac{\ln\!\left(f/f_1\right)}{\ln\!\left(f_2/f_1\right)},
  \label{eq:sweeptime}
\end{equation}
which is used to associate a time window with each excitation frequency.

Preprocessing consists of sorting the scan points along the cable axis,
removing dead channels (all-zero velocity), selecting the active shaker
channel by maximum energy, and correcting the known offset between the LDV
and shaker coordinate systems.

\paragraph{Band-pass filtering.}
% ldv_analysis/signal.py :: bandpass
Where a single excitation frequency $f_0$ is analysed, the velocities are
band-pass filtered with a zero-phase Butterworth filter of order 4
(applied forward and backward, \texttt{sosfiltfilt}, so the effective order
is 8 and the phase response is identically zero). The passband is
$f_0 \pm \SI{5}{Hz}$ (absolute bandwidth mode). Zero-phase filtering is
essential here: the coupling efficiency is a ratio of two simultaneously
filtered signals, and any filter-induced phase shift would bias it.

\paragraph{Integration to displacement.}
% ldv_analysis/signal.py :: integrate_fft
The LDV measures velocity; all coupling quantities are defined on
displacements. Integration is performed in the frequency domain. With
$V(f) = \mathcal{F}\{v\}(f)$ the discrete Fourier transform of the velocity,
the displacement spectrum follows from the differentiation property of the
Fourier transform,
\begin{equation}
  u(t) \;=\; \mathcal{F}^{-1}\!\left\{ \frac{V(f)}{i\,2\pi f} \right\}(t),
  \qquad f \ge f_\mathrm{min},
  \label{eq:integration}
\end{equation}
where the division is protected below $f_\mathrm{min} = \SI{0.5}{Hz}$
(frequencies below the floor are divided by $i\,2\pi f_\mathrm{min}$ instead)
to prevent the $1/f$ amplification of DC offset and low-frequency drift that
makes naive time-domain integration unstable. This is a linear, phase-exact
operation: a velocity component $v = \hat v\cos(2\pi f t)$ maps to
$u = \tfrac{\hat v}{2\pi f}\sin(2\pi f t)$, as expected.

\subsection{Geometry, Sag Determination and the Coupling Parameter $\Theta$}
\label{sec:proc:geometry}

\paragraph{Chord and axial projection.}
% ldv_analysis/geometry.py :: compute_chord, project_onto_chord
The cable spans a gap between two mounting points; the first and last scan
point inside the gap, $\mathbf{P}_1$ and $\mathbf{P}_2$, define the
\emph{chord}
\begin{equation}
  \hat{\mathbf{e}} \;=\; \frac{\mathbf{P}_2-\mathbf{P}_1}
                              {\lVert\mathbf{P}_2-\mathbf{P}_1\rVert},
  \qquad
  L \;=\; \lVert\mathbf{P}_2-\mathbf{P}_1\rVert .
  \label{eq:chord}
\end{equation}
Axial (in-line) displacement components are obtained by projection onto the
chord,
\begin{equation}
  u_\parallel(x_i, t) \;=\; \mathbf{u}(x_i,t)\cdot\hat{\mathbf{e}},
  \label{eq:projection}
\end{equation}
which is the displacement component a DAS fibre would sense. Projection onto
the chord (rather than taking the $x$-component) avoids mixing transverse
motion into the axial estimate when the cable is not perfectly aligned with
the coordinate axes.

\paragraph{Sag measurement.}
% ldv_analysis/geometry.py :: measure_sag (use_3d=True, use_parabola=True)
For every scan point in the gap, the perpendicular distance from the chord
line is computed in full 3-D:
\begin{equation}
  s_i = (\mathbf{P}_i - \mathbf{P}_1)\cdot\hat{\mathbf{e}},
  \qquad
  d_i = \bigl\lVert (\mathbf{P}_i-\mathbf{P}_1) - s_i\,\hat{\mathbf{e}}
        \bigr\rVert ,
  \label{eq:perpdist}
\end{equation}
where $s_i$ is the arc coordinate along the chord and $d_i$ the perpendicular
offset. Since the endpoints lie on the chord by construction
($d(0) = d(L) = 0$), a parabola \emph{constrained through both endpoints} is
fitted,
\begin{equation}
  d(s) \;\approx\; -a\, s\,(s - L),
  \qquad
  a \;=\; \frac{\sum_i \phi_i\, d_i}{\sum_i \phi_i^2},
  \quad \phi_i = s_i (s_i - L),
  \label{eq:parabolafit}
\end{equation}
a one-parameter linear least-squares problem. The midpoint sag is the vertex
value of the fitted parabola,
\begin{equation}
  w_0 \;=\; d(L/2) \;=\; \frac{a L^2}{4} \Big|_{a<0} ,
  \label{eq:sagvalue}
\end{equation}
and the root-mean-square residual of the fit,
$\sigma_{w_0} = \sqrt{\tfrac{1}{N}\sum_i (d_i - d(s_i))^2}$, is retained as
the sag measurement uncertainty (it captures outlier scan points and
deviations of the true shape from the parabolic approximation). The parabola
is the small-sag shape of a beam/cable under uniform load, so this fit is
consistent with the analytical model of
Section~\ref{sec:proc:geometry:theta}. If the fit yields $a \ge 0$ (no
physical downward arch), the maximum raw perpendicular distance is used
instead.

\paragraph{The coupling parameter $\Theta$.}
\label{sec:proc:geometry:theta}
% ldv_analysis/config.py :: theta_from_sag, theta_from_material
The bending-stress-relief model \citep{ProbstSimone} predicts that the
static strain-transfer efficiency of a suspended segment is
\begin{equation}
  \eta_\mathrm{stat} \;=\; \frac{\delta x_\ell}{\delta L}
                    \;=\; \frac{1}{1+\Theta},
  \qquad
  \Theta \;=\; \frac{EA}{L}\,\frac{(C_1 x_{f0})^2}{E I C_2},
  \label{eq:etatheta}
\end{equation}
where $x_{f0}$ is the gravity-induced midpoint sag,
$C_1 = \int_0^L (\phi')^2\,\mathrm{d}x$ and
$C_2 = \int_0^L (\phi'')^2\,\mathrm{d}x$ are the mode-shape constants of the
first symmetric clamped--clamped deflection shape
$\phi(x) = \tfrac{1}{2}\bigl[1-\cos(2\pi x/L)\bigr]$. Substituting the
tension-free equilibrium sag
\begin{equation}
  x_{f0} \;=\; w_0 \;=\; \frac{\rho A g L^4}{4\pi^4 E I}
  \label{eq:sagtheory}
\end{equation}
into Eq.~\eqref{eq:etatheta} gives the \emph{material form}
\begin{equation}
  \Theta_\mathrm{mat} \;=\; \frac{\rho^2 g^2 A^3 L^8}{128\,\pi^8 E^2 I^3},
  \label{eq:thetamaterial}
\end{equation}
while keeping the \emph{measured} sag $w_0$ gives the \emph{sag form}
\begin{equation}
  \Theta_\mathrm{sag} \;=\; \frac{A\, w_0^2}{8 I}
  \;\overset{\text{circular}}{=}\; \frac{1}{2}\left(\frac{w_0}{r}\right)^2 ,
  \label{eq:thetasag}
\end{equation}
where the last equality holds for a circular cross-section with
$A = \pi r^2$ and $I = \pi r^4/4$. For the rectangular ribbon cable
(width $w$, height $h$, bending about the weak axis) $A = wh$ and
$I = w h^3/12$ are used, so Eq.~\eqref{eq:thetasag} correctly generalises
beyond circular cables. The sag form is preferred throughout this work
because it absorbs the actual (pretension-affected) deployment geometry; the
material form is reported for comparison. Young's modulus enters as the mean
of three independent tensile tests per cable; $\Theta_\mathrm{mat}$ is
evaluated once per test value and reported as mean~$\pm$~standard deviation.

\subsection{Strain and Elongation Estimation}
\label{sec:proc:strain}

The coupling efficiency compares the \emph{axial elongation of the cable}
with the \emph{elongation imposed at its boundaries}. Three independent
estimators of the cable elongation are implemented --- M1 per-segment 3-D
arc-length, M2 spatial gradient and M3 Fourier-domain gradient --- and their
agreement is a built-in consistency check of the processing.

\subsubsection{Method 1: Per-Segment 3-D Arc-Length Elongation}
\label{sec:proc:strain:m1}
% ldv_analysis/strain.py :: strain_3d_arclength
Method 1 estimates the elongation of each segment between two adjacent scan
points directly from their 3-D displacement vectors, without projecting the
whole field onto the global chord first. Let
$\mathbf{r}_i = \mathbf{P}_{i+1} - \mathbf{P}_i$ be the static segment
vector, $L_{0,i} = \lVert\mathbf{r}_i\rVert$ its rest length and
$\hat{\mathbf{e}}_i = \mathbf{r}_i / L_{0,i}$ its \emph{local} unit vector.
The deformed segment length is
\begin{equation}
  L_i(t) \;=\; \bigl\lVert \mathbf{r}_i
      + \Delta\mathbf{u}_i(t) \bigr\rVert ,
  \qquad
  \Delta\mathbf{u}_i(t) = \mathbf{u}_{i+1}(t) - \mathbf{u}_i(t).
\end{equation}
Because the dynamic displacements ($\sim\si{\micro\meter}$) are many orders
of magnitude smaller than the segment lengths ($\sim\si{\milli\meter}$), the
square root is linearised by a first-order Taylor expansion:
\begin{equation}
  L_i(t) - L_{0,i}
  \;=\; L_{0,i}\sqrt{1
        + \frac{2\,\mathbf{r}_i\!\cdot\!\Delta\mathbf{u}_i}{L_{0,i}^2}
        + \frac{\lVert\Delta\mathbf{u}_i\rVert^2}{L_{0,i}^2}} - L_{0,i}
  \;\approx\;
  \hat{\mathbf{e}}_i \cdot \Delta\mathbf{u}_i(t)
  \;\equiv\; \delta d_i(t).
  \label{eq:strainM1}
\end{equation}
Geometrically: \emph{the relative displacement of the two points is projected
onto the local segment direction} $\hat{\mathbf{e}}_i$ — only the component
of relative motion along the segment changes its length to first order; the
perpendicular component merely rotates it (a second-order length effect that
the linearisation deliberately drops, consistent with the small-strain
regime). For a sagged cable the $\hat{\mathbf{e}}_i$ follow the local cable
direction, so Method 1 measures the \emph{arc-length} change, which is
exactly the quantity a strained optical fibre integrates. The local strain
estimate of segment $i$ is $\varepsilon_i = \delta d_i / L_{0,i}$, and the
total cable elongation is the telescoping sum over the gap,
\begin{equation}
  \delta x_\ell(t) \;=\; \sum_{i\,\in\,\mathrm{gap}} \delta d_i(t).
  \label{eq:deltaxl}
\end{equation}
Method 1 is the primary estimator in this work: it makes no assumption
about the smoothness of the displacement field between scan points and is
insensitive to the chord-projection error for sagged geometries.

\subsubsection{Method 2: Spatial Gradient of the Axial Displacement}
% ldv_analysis/strain.py :: strain_spatial_gradient
Axial strain is the spatial derivative of the axial displacement along the
cable,
\begin{equation}
  \varepsilon(s,t) \;=\; \frac{\partial\, u_\parallel(s,t)}{\partial s},
  \label{eq:strainM2def}
\end{equation}
evaluated with second-order central finite differences on the non-uniform
scan-point grid. For interior point $i$ with spacings
$\Delta s_- = s_i - s_{i-1}$ and $\Delta s_+ = s_{i+1} - s_i$:
\begin{equation}
  \varepsilon_i \;=\;
    u_{\parallel,i+1}\,\frac{\Delta s_-}{\Delta s_+(\Delta s_+ + \Delta s_-)}
  + u_{\parallel,i}\,\frac{\Delta s_+ - \Delta s_-}{\Delta s_+ \Delta s_-}
  - u_{\parallel,i-1}\,\frac{\Delta s_+}{\Delta s_-(\Delta s_+ + \Delta s_-)} .
  \label{eq:strainM2}
\end{equation}
This reduces to the familiar $(u_{i+1}-u_{i-1})/(2\Delta s)$ for uniform
spacing and is exact for locally quadratic displacement fields. The total
elongation follows by trapezoidal integration over the gap,
$\delta x_\ell^{(2)}(t) = \int_\mathrm{gap} \varepsilon(s,t)\,\mathrm{d}s$.

\subsubsection{Method 3: Fourier-Domain Spatial Gradient}
% ldv_analysis/strain.py :: strain_fourier_gradient
As a cross-check with different numerical error characteristics, the axial
displacement is resampled onto a uniform grid $s_k$ (cubic interpolation)
and differentiated spectrally using the Fourier differentiation theorem,
\begin{equation}
  \varepsilon(s,t) \;=\;
  \mathcal{F}_s^{-1}\bigl\{\, i k\, \mathcal{F}_s\{u_\parallel\}\,\bigr\},
  \label{eq:strainM3}
\end{equation}
with $k$ the spatial angular wavenumber. Spectral differentiation is exact
for band-limited periodic signals; since the cable displacement is not
periodic over the gap, Gibbs-type artefacts can appear near the endpoints,
which is why only gap-interior stations enter the elongation integral. In
practice Method 3 agrees with Methods 1--2 in the smooth interior and serves
to expose discretisation artefacts of the finite-difference estimate.

\subsubsection{Reference Elongations}
% ldv_analysis/freqdomain.py :: build_coupling_signals (x_source)
Three references for the imposed boundary elongation $\delta L(t)$ are
available:
\begin{itemize}
  \item \textbf{Endpoint pair} (default): the span change measured on the
    cable itself,
    $\delta L = u_\parallel(\mathbf{P}_2) - u_\parallel(\mathbf{P}_1)$.
    It shares its endpoints with $\delta x_\ell$, so the efficiency isolates
    bending stress relief and is immune to cross-instrument calibration.
  \item \textbf{Shaker}: the chord-projected shaker-head displacement
    $\delta L = \mathbf{u}_\mathrm{sh}\cdot\hat{\mathbf{e}}$ (sign-corrected
    so that positive means stretching), i.e.\ the imposed ground motion.
  \item \textbf{Shaker-to-end}: the elongation of the sub-span between the
    shaker and the adjacent endpoint sensor, used to diagnose losses at the
    first contact.
\end{itemize}

\subsection{Time-Domain Coupling Efficiency (Quality Control)}
\label{sec:proc:timedomain}
% ldv_analysis/analysis.py :: per_frequency_qc

For inspection at a single excitation frequency $f_0$, the record is windowed
to $t_c(f_0) \pm n_\mathrm{cyc}/f_0$ (Eq.~\ref{eq:sweeptime},
$n_\mathrm{cyc}=5$), band-passed (Section~\ref{sec:proc:preprocessing}),
integrated to displacement, and the instantaneous efficiency is formed:
\begin{equation}
  \eta(t) \;=\; \frac{\delta x_\ell(t)}{\delta L(t)},
  \qquad
  \eta_\mathrm{med} = \operatorname{median}_t\, \eta(t)
  \quad\text{over } |\delta L(t)| > 0.05\,\max|\delta L| .
  \label{eq:etatime}
\end{equation}
The amplitude gate excludes the zero crossings of the reference, where the
ratio is ill-conditioned; the median suppresses the remaining outliers. A
phase-insensitive variant replaces both signals by their analytic-signal
envelopes (Section~\ref{sec:proc:amplitude}),
$\eta_\mathrm{env}(t) = |\mathcal{H}\{\delta x_\ell\}| /
|\mathcal{H}\{\delta L\}|$: comparing $\eta_\mathrm{med}$ and
$\eta_\mathrm{env}$ distinguishes whether a low efficiency is caused by an
amplitude deficit or by a phase lag between cable and reference.

These time-domain values are used for quality control and for illustrating
the mechanism at single frequencies; all summary efficiencies in the results
chapter are computed in the frequency domain
(Section~\ref{sec:proc:freqdomain}), which uses the whole record at once and
is considerably more stable.

\subsection{Frequency-Domain Transfer Functions}
\label{sec:proc:freqdomain}
% ldv_analysis/freqdomain.py

\subsubsection{Signals}
From the full (unfiltered) sweep record, integrated to displacement, three
signals are built:
\begin{align}
  X(t) &= \delta L(t) && \text{input: boundary span change,} \notag\\
  Y_1(t) &= \delta x_\ell(t) && \text{output 1: cable elongation
                                (Eq.~\ref{eq:deltaxl}),} \notag\\
  Y_2(t) &= u_{\perp,\mathrm{mid}}(t) && \text{output 2: midpoint deflection
                                along the static sag direction.} \notag
\end{align}
$H_\eta(f) = Y_1/X$ is the dynamic strain-transfer efficiency;
$H_\mathrm{mid}(f) = Y_2/X$ measures the intensity of the bending mode, whose
resonance structure explains the frequency dependence of $H_\eta$
\citep{ProbstSimone}.

\subsubsection{Welch Spectral Estimation}
% ldv_analysis/freqdomain.py :: welch_spectra
The record is divided into $K$ segments of length
$n_\mathrm{seg} = 2 f_s = 10\,000$ samples (\SI{2}{s}, frequency resolution
$\Delta f = \SI{0.5}{Hz}$) with 50\,\% overlap, each tapered with a Hann
window $w[n]$. With $X_k(f)$, $Y_k(f)$ the windowed DFTs of segment $k$,
the auto- and cross-spectral densities are the segment averages
\begin{equation}
  S_{XX}(f) = \frac{1}{K}\sum_{k=1}^{K} |X_k(f)|^2 , \qquad
  S_{YY}(f) = \frac{1}{K}\sum_{k=1}^{K} |Y_k(f)|^2 , \qquad
  S_{XY}(f) = \frac{1}{K}\sum_{k=1}^{K} X_k^*(f)\, Y_k(f),
  \label{eq:welch}
\end{equation}
(up to the standard window normalisation, which cancels in every ratio used
below). The averaging over $K \approx 11$ segments reduces the variance of
each spectral estimate by $\approx 1/K$ — this is what makes the Welch-based
transfer functions far smoother than any single-FFT estimate. The
\emph{ordinary coherence}
\begin{equation}
  \gamma^2(f) \;=\; \frac{|S_{XY}(f)|^2}{S_{XX}(f)\, S_{YY}(f)}
  \;\in\; [0, 1]
  \label{eq:coherence}
\end{equation}
measures the fraction of output power linearly explained by the input;
$\gamma^2 < 1$ indicates noise, nonlinearity, or unmeasured inputs.

\subsubsection{FRF Estimators: Direct, $H_1$, $H_2$, $H_v$}
\label{sec:proc:estimators}
% ldv_analysis/freqdomain.py :: direct_transfer, estimate_H

\paragraph{Direct estimate.}
The conceptually simplest frequency response estimate is the ratio of the
full-record windowed Fourier transforms,
\begin{equation}
  H_\mathrm{direct}(f) \;=\;
  \frac{\mathcal{F}\{w\,Y\}(f)}{\mathcal{F}\{w\,X\}(f)} ,
  \label{eq:Hdirect}
\end{equation}
evaluated on the dense grid $\Delta f_\mathrm{d} = f_s/N = \SI{0.083}{Hz}$.
Every frequency bin is a single realisation: the estimate is unbiased in the
noise-free limit but has no variance reduction, so amplitude and phase
scatter strongly wherever excitation or response is weak.

Comparing Eq.~\eqref{eq:Hdirect} directly against $H_1$ would conflate two
distinct effects — the estimator algebra and the amount of spectral averaging.
To separate them, the direct spectra are smoothed along the \emph{frequency}
axis (Daniell smoothing) rather than over time segments: with a normalised
kernel $k$ of $m$ bins,
\begin{equation}
  \tilde S_{AB}(f_i) = \sum_{j} k_j\, A^*(f_{i+j})\, B(f_{i+j}),
  \qquad
  H_\mathrm{direct}^{\mathrm{sm}}(f) = \frac{\tilde S_{XY}(f)}{\tilde S_{XX}(f)} .
  \label{eq:Hdirectsmooth}
\end{equation}
Smoothing the products \emph{before} the division makes the estimate
input-power weighted, so bins with weak $|X|$ contribute little instead of
diverging; this also removes the positive bias that the unweighted ratio
carries at low SNR. A kernel with $\sum_j k_j = 1$ averages
$n_\mathrm{eff} = 1/\sum_j k_j^2$ independent bins, which sets both the
variance reduction and the effective resolution bandwidth. The kernel width is
therefore chosen so that $n_\mathrm{eff}\,\Delta f_\mathrm{d}$ equals the
\emph{equivalent noise bandwidth} of the Welch estimate,
\begin{equation}
  B_\mathrm{ENBW} = f_s \frac{\sum_n w[n]^2}{\left(\sum_n w[n]\right)^2}
  = 1.5\,\Delta f = \SI{0.75}{Hz}
  \quad\text{(Hann)},
  \label{eq:enbw}
\end{equation}
i.e.\ $n_\mathrm{eff} \approx 9$ dense bins. Note that
$B_\mathrm{ENBW} > \Delta f$: the Hann taper spreads each spectral line over
1.5 bins, so a Welch estimate does not resolve to its nominal bin spacing, and
matching $B_\mathrm{ENBW}$ rather than $\Delta f$ is what puts the two
estimates on an equal footing. With this choice the smoothed direct estimate
and $H_1$ agree to within a few percent in $|H|$ while retaining the dense
frequency grid; the residual difference is attributable to the estimator
algebra alone. Both the raw and the smoothed direct curves are shown, so the
figure separates the effect of averaging (raw $\to$ smoothed) from the effect
of the estimator (smoothed $\to$ $H_1$).

\paragraph{$H_1$ — noise on the output.}
Assume the measured output contains additive noise uncorrelated with the
input, $Y = H X + N$. Minimising the output residual over the segment
ensemble,
$\min_H \sum_k |Y_k - H X_k|^2$, gives the least-squares solution
\begin{equation}
  H_1(f) \;=\; \frac{S_{XY}(f)}{S_{XX}(f)} .
  \label{eq:H1}
\end{equation}
Since $\mathbb{E}\{X^* N\} = 0$, the cross-spectrum averages the output noise
away and $H_1$ is unbiased under output noise. If the \emph{input} channel is
noisy, $S_{XX}$ is inflated and $|H_1|$ \emph{under}estimates the true FRF.

\paragraph{$H_2$ — noise on the input.}
Assume instead noise on the input, $X = X_0 + M$ with $Y = H X_0$.
Minimising the input-referenced residual
$\min_H \sum_k |X_k - Y_k / H|^2$ yields
\begin{equation}
  H_2(f) \;=\; \frac{S_{YY}(f)}{S_{YX}(f)}
        \;=\; \frac{S_{YY}(f)}{S_{XY}^*(f)} .
  \label{eq:H2}
\end{equation}
$H_2$ is unbiased under input noise but \emph{over}estimates $|H|$ when the
output is noisy. The two estimators bracket the true FRF,
\begin{equation}
  |H_1(f)| \;\le\; |H(f)| \;\le\; |H_2(f)|,
  \qquad
  \frac{|H_1(f)|}{|H_2(f)|} = \gamma^2(f),
  \label{eq:bracket}
\end{equation}
so their separation is itself a coherence display: the estimators agree
exactly where $\gamma^2 = 1$ and diverge where the data are noisy.

\paragraph{$H_v$ — noise on both channels (total least squares).}
When both channels carry comparable noise — realistic here, since input and
output are both derived from LDV velocities through the same processing —
the total-least-squares estimator distributes the residual over both
channels. Writing the spectral sample covariance of the vector $(X, Y)$ per
frequency bin,
\begin{equation}
  \mathbf{C}(f) \;=\;
  \begin{pmatrix} S_{XX} & S_{XY}^* \\ S_{XY} & S_{YY} \end{pmatrix},
\end{equation}
the TLS solution is the direction of the eigenvector belonging to the
\emph{smallest} eigenvalue $\lambda_-$ of $\mathbf{C}$ (the noise
eigenvalue). For the $2\times 2$ Hermitian matrix this eigenvalue is
\begin{equation}
  \lambda_-(f) \;=\; \frac{S_{XX}+S_{YY}}{2}
    - \sqrt{\left(\frac{S_{XX}-S_{YY}}{2}\right)^{\!2} + |S_{XY}|^2},
\end{equation}
and solving the eigenvector relation for the FRF gives the closed form
\begin{equation}
  H_v(f) \;=\; \frac{S_{YY}(f) - \lambda_-(f)}{S_{XY}^*(f)} .
  \label{eq:Hv}
\end{equation}
One verifies directly that in the noise-free case ($\lambda_- = 0$,
$S_{YY} = |H|^2 S_{XX}$, $S_{XY} = H S_{XX}$) Eq.~\eqref{eq:Hv} returns
exactly $H$, and that $|H_1| \le |H_v| \le |H_2|$ holds bin-by-bin. All
three Welch estimators share the same phase,
$\arg H_1 = \arg H_2 = \arg H_v = \arg S_{XY}$, so the estimator choice
affects only the magnitude.

\paragraph{Phase presentation.}
Phases are reported \emph{wrapped} to $(-180^\circ, 180^\circ]$. Unwrapped
phase curves accumulate spurious $\pm 360^\circ$ jumps whenever
low-coherence bins randomise the phase, which makes them diverge without
physical meaning; the wrapped display is immune to this and the physically
relevant transition (the $\approx 180^\circ$ flip across the bending
resonance) remains fully visible.

\subsubsection{Band-Averaged Efficiency and Coherence Gating}
% ldv_analysis/freqdomain.py :: band_mean_eta
The quasi-static efficiency reported per configuration is the coherence-gated
band average
\begin{equation}
  \bar\eta \;=\; \bigl\langle\, |H_\eta(f)| \,\bigr\rangle_{f \in B},
  \qquad
  B = \{ f \in [f_\mathrm{min}, f_\mathrm{max}] :
         \gamma^2_\eta(f) \ge 0.7 \},
  \label{eq:bandeta}
\end{equation}
with $[f_\mathrm{min}, f_\mathrm{max}] = [5, 25]\,\si{Hz}$ well below the
first resonance. The gate removes noise-dominated bins that would otherwise
bias the average; the standard deviation over $B$ is reported as the error
bar.

\subsubsection{Single-Frequency Cross-Check and Resonance Characterisation}
% ldv_analysis/freqdomain.py :: stft_point_transfer, fit_lorentzian
As an independent validation, point estimates
$H(f_0) = \mathcal{F}\{w\,Y\}(f_0)\,/\,\mathcal{F}\{w\,X\}(f_0)$ are computed
from short Hann-tapered windows centred on the sweep arrival $t_c(f_0)$ and
overlaid on the Welch curves.

The fundamental resonance $f_1$ is identified as the global maximum of
$|H_\mathrm{mid}(f)|$; a driven single-degree-of-freedom magnitude response
\begin{equation}
  |H(f)| \;=\; \frac{G\, f_1^2}
   {\sqrt{\left(f_1^2 - f^2\right)^2 + \left(f_1 f / Q\right)^2}}
  \label{eq:sdof}
\end{equation}
is fitted in a window around the peak, yielding the quality factor $Q$ and
damping ratio $\zeta = 1/(2Q)$. The theoretical clamped--clamped fundamental
\begin{equation}
  f_1^\mathrm{th} \;=\; \frac{\alpha_1^2}{2\pi L^2}\sqrt{\frac{EI}{\rho A}},
  \qquad \alpha_1^2 = 4.730^2 \approx 22.37,
  \label{eq:f1theory}
\end{equation}
is quoted for comparison (cf.~Section~\ref{sec:proc:modal}).

\subsection{Amplitude Extraction}
\label{sec:proc:amplitude}
% ldv_analysis/amplitude.py

Beyond ratios, the absolute displacement and strain amplitudes are needed as
reference values for the numerical simulations. Two extraction routes are
implemented and cross-validated.

\paragraph{Time domain: analytic-signal envelope.}
For a windowed, band-passed trace $x(t)$ the analytic signal is
$x_a(t) = x(t) + i\,\mathcal{H}\{x\}(t)$ with $\mathcal{H}$ the Hilbert
transform; its magnitude $|x_a(t)|$ is the instantaneous envelope. The
representative amplitude is the envelope median (robust against the
filter-edge droop at the window boundaries); envelope maximum and
$\sqrt{2}\,\mathrm{RMS}$ variants are available for comparison.

\paragraph{Frequency domain: band-energy read-out.}
The velocities are integrated once per dataset; at each target frequency
$f_0$ the displacement signal is windowed around $t_c(f_0)$, Hann-tapered
and Fourier transformed. Because the log-sweep smears the tone over several
bins (especially at low frequency, where the window spans many seconds), the
amplitude is recovered from the \emph{energy in the analysis band}
$B = [f_0 - 5, f_0 + 5]\,\si{Hz}$ rather than from a single bin. For a tone
of peak amplitude $A$ under a window $w[n]$, Parseval's theorem gives
\begin{equation}
  A \;=\; 2\sqrt{\frac{\sum_{k \in B} |X_k|^2}{N \sum_n w[n]^2}} ,
  \label{eq:bandenergy}
\end{equation}
which is exact for an isolated in-band tone. This route avoids the
per-frequency band-pass filtering of the full record and is roughly two
orders of magnitude faster, enabling dense frequency grids; it agrees with
the envelope route to within a few percent below $\approx\SI{100}{Hz}$.

\subsection{Extension--Compression Asymmetry}
\label{sec:proc:asymmetry}
% ldv_analysis/asymmetry.py

A sagged segment stiffens when its endpoints are pulled apart (the sag is
straightened) and softens when they are pushed together (the sag deepens);
numerical simulations show correspondingly different responses for the two
deformation directions. A perfectly linear system responds
mirror-symmetrically to the two half-cycles of a sinusoidal drive, so any
measurable asymmetry is a signature of this geometric nonlinearity. Three
complementary detectors are evaluated per excitation frequency:

\begin{enumerate}
  \item \textbf{Half-cycle conditional efficiency.} The samples of the
  reference elongation are split by sign — extension
  ($\delta L > 0$, endpoints separating) versus compression
  ($\delta L < 0$) — with a $25\,\%$-of-maximum amplitude gate excluding the
  zero crossings. The efficiencies
  \begin{equation}
    \eta_\mathrm{ext} = \operatorname*{median}_{\delta L(t) > \delta_g}
      \frac{\delta x_\ell(t)}{\delta L(t)},
    \qquad
    \eta_\mathrm{comp} = \operatorname*{median}_{\delta L(t) < -\delta_g}
      \frac{\delta x_\ell(t)}{\delta L(t)}
    \label{eq:etahalfcycle}
  \end{equation}
  coincide for a linear system; $\eta_\mathrm{ext} \ne \eta_\mathrm{comp}$
  quantifies the directional coupling difference.

  \item \textbf{Peak asymmetry.} From the mean positive-peak amplitude
  $A_+$ and mean negative-peak magnitude $A_-$ of the cable elongation
  $\delta x_\ell(t)$:
  \begin{equation}
    \alpha \;=\; \frac{A_+ - A_-}{A_+ + A_-} \;\in\; [-1, 1],
    \label{eq:peakasym}
  \end{equation}
  with $\alpha > 0$ meaning the cable elongates further during extension
  half-cycles than it shortens during compression.

  \item \textbf{Even-harmonic distortion.} Expanding the response in powers
  of the drive, $y = c_1 x + c_2 x^2 + c_3 x^3 + \dots$ with
  $x = A\cos\omega t$, the quadratic term produces
  $c_2 A^2 \cos^2\omega t = \tfrac{c_2 A^2}{2}(1 + \cos 2\omega t)$ — energy
  at \emph{twice} the drive frequency — whereas the cubic (symmetric) term
  appears at $3\omega$. Since a sign-asymmetric stiffness is exactly a
  quadratic term, the ratio
  \begin{equation}
    R_2 \;=\; \frac{|Y(2 f_0)|}{|Y(f_0)|}
    \label{eq:h2ratio}
  \end{equation}
  isolates the asymmetric part of the response independently of the
  reference channel; $R_3 = |Y(3f_0)|/|Y(f_0)|$ is reported alongside to
  attribute nonlinearity to the symmetric (softening/stiffening) branch.
\end{enumerate}

For the linearity data set (Section~\ref{sec:proc:linearity}) the half-cycle
efficiencies are additionally evaluated in bins of the instantaneous drive
amplitude, revealing whether an asymmetry develops progressively with
amplitude (e.g.\ slack forming on the compression half-cycle only).

\subsection{Modal Analysis}
\label{sec:proc:modal}
% ldv_analysis/modal.py

\subsubsection{Detection Fingerprint and Spectral Smoothing}
% modal.spectral_fingerprint
Transverse resonances are detected on a \emph{fingerprint}: per-sensor
spectra of the transverse velocities $v_y, v_z$ over the gap, collapsed
across sensors into the total-power detection signal
\begin{equation}
  P(f) \;=\; \sum_{s\,\in\,\mathrm{gap}}
      \Bigl( |V_y(f; s)|^2 + |V_z(f; s)|^2 \Bigr).
  \label{eq:fingerprint}
\end{equation}
The per-sensor spectra are Welch estimates
($n_\mathrm{seg} = 2 f_s$, $\Delta f = \SI{0.5}{Hz}$, Hann, 50\,\% overlap,
$K \approx 11$ averages) rather than raw single-record FFTs: the
$\mathrm{var} \propto 1/K$ averaging suppresses the bin-to-bin fluctuations
that previously caused the peak picker to fire on noise wiggles. An optional
moving-average (boxcar) smoothing over a configurable bandwidth can be
applied on top. Resonances are then the peaks of $P(f)$ (normalised to its
in-band maximum) exceeding a prominence threshold with a minimum mutual
spacing (\texttt{scipy.signal.find\_peaks}).

\paragraph{Input-bias correction.}
% modal.shaker_input_spectrum, normalize_by_input
The shaker head does not excite purely axially: its own recording shows
frequency-dependent energy in the transverse components. A spectral peak in
the cable response may therefore reflect a peak in the \emph{excitation}
rather than a cable resonance. Two diagnostics address this. First, the
input directionality is quantified by the per-frequency power fractions of
the shaker-head components,
\begin{equation}
  \phi_c(f) \;=\; \frac{|V_c^\mathrm{sh}(f)|^2}
   {\sum_{c' \in \{x,y,z\}} |V_{c'}^\mathrm{sh}(f)|^2},
  \qquad c \in \{x, y, z\}.
  \label{eq:directionality}
\end{equation}
Second, the detection signal can be normalised by the total shaker input
power, $P_\mathrm{norm}(f) = P(f) / P_\mathrm{sh}(f)$ — an FRF-like quantity
in which input coloration cancels: peaks that survive the normalisation are
genuine output-per-input amplifications (resonances), while peaks that
disappear were inherited from the excitation spectrum.

\subsubsection{Complex Mode-Shape Extraction}
% modal.extract_mode_shape
For each detected resonance $f_n$, the transverse velocities are band-passed
to $f_n(1 \pm 0.1)$, integrated to displacement, and the complex Fourier
coefficient at the bin nearest $f_n$ is read per sensor:
\begin{equation}
  \varphi_y(s) = U_y(f_n; s), \qquad \varphi_z(s) = U_z(f_n; s).
  \label{eq:modeshape}
\end{equation}
The complex-valued shapes retain the relative phase between sensors, so
travelling-wave content and sign flips across nodes are preserved.

\subsubsection{Euler--Bernoulli Reference Modes}
% modal.clamped_clamped_mode
The gap is modelled as a doubly-clamped uniform beam. Free vibration of an
Euler--Bernoulli beam obeys $EI\, w'''' = \rho A\, \omega^2 w$; with clamped
boundary conditions $w(0)=w'(0)=w(L)=w'(L)=0$, non-trivial solutions exist
only at the roots $\alpha_n$ of the frequency equation
\begin{equation}
  \cos\alpha \cosh\alpha \;=\; 1,
  \qquad
  \alpha_n = 4.730,\ 7.853,\ 10.996,\ 14.137,\ 17.279,\ \dots
  \label{eq:freqequation}
\end{equation}
with natural frequencies and mode shapes
\begin{equation}
  f_n \;=\; \frac{\alpha_n^2}{2\pi L^2} \sqrt{\frac{EI}{\rho A}},
  \qquad
  \varphi_n(\xi) = \bigl[\cosh\alpha_n\xi - \cos\alpha_n\xi\bigr]
   - \sigma_n \bigl[\sinh\alpha_n\xi - \sin\alpha_n\xi\bigr],
  \label{eq:ebmodes}
\end{equation}
where $\xi = x/L \in [0,1]$ and
$\sigma_n = (\cosh\alpha_n - \cos\alpha_n)/(\sinh\alpha_n - \sin\alpha_n)$.
The frequency ratios $f_n / f_1 = (\alpha_n/\alpha_1)^2 =
1,\ 2.76,\ 5.40,\ 8.93,\ \dots$ provide a material-free consistency check of
the mode ordering.

\subsubsection{Modal Assurance Criterion}
% modal.mac
Each measured shape is assigned to a theoretical mode by the Modal Assurance
Criterion,
\begin{equation}
  \operatorname{MAC}(\boldsymbol{\varphi}_m, \boldsymbol{\varphi}_t)
  \;=\;
  \frac{\bigl|\boldsymbol{\varphi}_m^{\mathsf{H}}
              \boldsymbol{\varphi}_t\bigr|^2}
       {\bigl(\boldsymbol{\varphi}_m^{\mathsf{H}}\boldsymbol{\varphi}_m\bigr)
        \bigl(\boldsymbol{\varphi}_t^{\mathsf{H}}\boldsymbol{\varphi}_t\bigr)}
  \;\in\; [0, 1],
  \label{eq:mac}
\end{equation}
where $\boldsymbol{\varphi}_m$ is the measured complex shape (the dominant
polarization component, sampled at the sensor positions),
$\boldsymbol{\varphi}_t$ the theoretical shape $\varphi_n(\xi_i)$ evaluated
at the same positions, and $^{\mathsf{H}}$ the conjugate transpose. The MAC
is the squared modulus of the normalised inner product — a
\emph{shape-correlation} coefficient: it equals 1 when the vectors are
collinear (identical shape up to any complex scale), and 0 when they are
orthogonal. Because a global complex factor cancels, the measured shape
needs no phase alignment or amplitude scaling before comparison, and a
$180^\circ$ phase flip across a node is handled correctly. The measured
resonance is assigned the mode order with the largest MAC; low best-MAC
values flag shapes that no beam mode explains (e.g.\ rigid-body-like motion
from mount compliance or shaker cross-talk).

\subsubsection{Polarization Analysis}
% modal.polarization_ellipse
At each sensor the pair $(\varphi_y, \varphi_z)$ defines an ellipse traced
by $\Re\{(\varphi_y, \varphi_z)\, e^{i\omega t}\}$ in the $y$--$z$ plane.
Decomposing into counter-rotating circular components
\begin{equation}
  A_+ = \tfrac{1}{2}\bigl(\varphi_y + i\,\varphi_z\bigr),
  \qquad
  A_- = \tfrac{1}{2}\bigl(\varphi_y^* + i\,\varphi_z^*\bigr)
  \quad\Longrightarrow\quad
  \text{major} = |A_+| + |A_-|,\quad
  \text{minor} = \bigl||A_+| - |A_-|\bigr|,
  \label{eq:ellipse}
\end{equation}
with axis ratio $\mathrm{minor}/\mathrm{major}$ (0 = linear, 1 = circular),
orientation $\theta = \tfrac{1}{2}(\arg A_+ + \arg A_-)$ and handedness
$\operatorname{sign}\Im(\varphi_y \varphi_z^*)$. A resonance is classified
linear / elliptical / circular by the amplitude-weighted majority over the
significant sensors.

\subsection{Linearity Analysis}
\label{sec:proc:linearity}
% ldv_analysis/analysis.py :: linearity_qc

A dedicated test signal probes amplitude dependence: eight discrete
frequencies (5--400\,\si{Hz}), each driven for \SI{3.5}{s} with a linear
amplitude ramp $0 \to 1$ followed by a \SI{0.5}{s} fade. Each segment is
extracted in full, band-passed ($f_0 \pm \SI{5}{Hz}$), integrated, and the
efficiency $\eta(t)$ (Eq.~\ref{eq:etatime}) is evaluated against the
\emph{instantaneous drive amplitude}, defined as the analytic-signal envelope
$a(t) = |\mathcal{H}\{\delta L\}(t)|$. Because the ramp sweeps amplitude
monotonically, plotting the binned median of $\eta$ against $a$ directly
tests linearity: a flat curve means amplitude-independent coupling, a slope
or hysteresis loop indicates nonlinearity (e.g.\ frictional slip at the
mounts above a threshold amplitude). A complementary slip indicator divides
the envelope of a single cable-sensor velocity by the envelope of the shaker
reference; for slip-free coupling this ratio stays constant along the ramp.
The frequency-domain cross-check estimates $H_1$
(Eq.~\ref{eq:H1}) on each mono-frequency segment and compares
$|H_1(f_0)|$ with the time-domain median efficiency.

\subsection{Uncertainty Treatment}
\label{sec:proc:uncertainty}

\begin{itemize}
  \item \textbf{Young's modulus:} three tensile tests per cable;
    $\Theta_\mathrm{mat}$ is evaluated per test and reported as mean $\pm$
    standard deviation. Since $\Theta_\mathrm{mat} \propto E^{-2}$, a
    relative modulus error $\epsilon_E$ propagates as
    $2\,\epsilon_E$ into $\Theta$.
  \item \textbf{Gap length:} the mounting-gap reading error is
    $\sigma_L = \SI{0.5}{cm}$. The material form is extremely sensitive,
    $\Theta_\mathrm{mat} \propto L^8 \Rightarrow
    \sigma_\Theta/\Theta = 8\,\sigma_L/L$ (40\,\% at $L=\SI{10}{cm}$!),
    which is the quantitative reason the sag-based
    $\Theta_\mathrm{sag}$ — independent of $L$ — is preferred.
  \item \textbf{Sag:} the RMS residual of the constrained parabola fit
    (Eq.~\ref{eq:parabolafit}) serves as $\sigma_{w_0}$; with
    $\Theta_\mathrm{sag} \propto w_0^2$,
    $\sigma_\Theta/\Theta = 2\,\sigma_{w_0}/w_0$.
  \item \textbf{Efficiency:} the standard deviation of $|H_\eta|$ over the
    coherence-gated quasi-static band (Eq.~\ref{eq:bandeta}); the
    $H_1$/$H_2$ bracket (Eq.~\ref{eq:bracket}) provides an additional
    systematic error envelope.
\end{itemize}

% ─── Bibliography hooks ─────────────────────────────────────────────────────
% \citep{ProbstSimone}  : S. Probst et al., "Unburied DAS on the Moon:
%                         Modeling Ground-to-Cable Strain Transfer" (draft).
% Probst et al. 2026    : Controlled Source DAS Coupling Tests, Earth and
%                         Space Science 13(3), e2025EA004817.
% Welch 1967            : P. Welch, IEEE Trans. Audio Electroacoust. 15(2).
% Allemang 2003         : R. Allemang, "The Modal Assurance Criterion --
%                         Twenty Years of Use and Abuse", Sound & Vibration.
% Bendat & Piersol 2010 : Random Data: Analysis and Measurement Procedures
%                         (H1/H2/Hv estimators, coherence).

% Required packages: amsmath, amssymb, siunitx, booktabs (for the tables).
% Every equation cites the implementing function in ldv_analysis/ in a margin
% comment so the text can be checked against the code.
% ═══════════════════════════════════════════════════════════════════════════

\section{Signal Processing and Data Analysis}
\label{sec:processing}

All quantitative results in this thesis are derived from the raw Laser
Doppler Vibrometer (LDV) recordings through the processing chain described in
this chapter. The chain is implemented in the Python package
\texttt{ldv\_analysis}; each subsection below corresponds to one module of
that package, and equations are stated exactly as implemented.

\subsection{Measurement Data and Preprocessing}
\label{sec:proc:preprocessing}

Each measurement consists of two synchronous recordings: the scanning LDV
measures the three-component velocity field
$\mathbf{v}(x_i, t) = (v_x, v_y, v_z)$ at $N_s$ scan points along the cable,
and a reference channel records the three-component velocity of the shaker
head, $\mathbf{v}_\mathrm{sh}(t)$. The excitation is a logarithmic sweep from
$f_1 = \SI{1}{Hz}$ to $f_2 = \SI{500}{Hz}$ over $T = \SI{10}{s}$ (after a
\SI{1}{s} buffer), sampled at $f_s = \SI{5}{kHz}$ with $N = 60\,000$ samples.
The instantaneous sweep frequency reaches $f$ at
% ldv_analysis/config.py :: sweep_time_of_frequency
\begin{equation}
  t_c(f) \;=\; t_\mathrm{buf} + T\,
    \frac{\ln\!\left(f/f_1\right)}{\ln\!\left(f_2/f_1\right)},
  \label{eq:sweeptime}
\end{equation}
which is used to associate a time window with each excitation frequency.

Preprocessing consists of sorting the scan points along the cable axis,
removing dead channels (all-zero velocity), selecting the active shaker
channel by maximum energy, and correcting the known offset between the LDV
and shaker coordinate systems.

\paragraph{Band-pass filtering.}
% ldv_analysis/signal.py :: bandpass
Where a single excitation frequency $f_0$ is analysed, the velocities are
band-pass filtered with a zero-phase Butterworth filter of order 4
(applied forward and backward, \texttt{sosfiltfilt}, so the effective order
is 8 and the phase response is identically zero). The passband is
$f_0 \pm \SI{5}{Hz}$ (absolute bandwidth mode). Zero-phase filtering is
essential here: the coupling efficiency is a ratio of two simultaneously
filtered signals, and any filter-induced phase shift would bias it.

\paragraph{Integration to displacement.}
% ldv_analysis/signal.py :: integrate_fft
The LDV measures velocity; all coupling quantities are defined on
displacements. Integration is performed in the frequency domain. With
$V(f) = \mathcal{F}\{v\}(f)$ the discrete Fourier transform of the velocity,
the displacement spectrum follows from the differentiation property of the
Fourier transform,
\begin{equation}
  u(t) \;=\; \mathcal{F}^{-1}\!\left\{ \frac{V(f)}{i\,2\pi f} \right\}(t),
  \qquad f \ge f_\mathrm{min},
  \label{eq:integration}
\end{equation}
where the division is protected below $f_\mathrm{min} = \SI{0.5}{Hz}$
(frequencies below the floor are divided by $i\,2\pi f_\mathrm{min}$ instead)
to prevent the $1/f$ amplification of DC offset and low-frequency drift that
makes naive time-domain integration unstable. This is a linear, phase-exact
operation: a velocity component $v = \hat v\cos(2\pi f t)$ maps to
$u = \tfrac{\hat v}{2\pi f}\sin(2\pi f t)$, as expected.

\subsection{Geometry, Sag Determination and the Coupling Parameter $\Theta$}
\label{sec:proc:geometry}

\paragraph{Chord and axial projection.}
% ldv_analysis/geometry.py :: compute_chord, project_onto_chord
The cable spans a gap between two mounting points; the first and last scan
point inside the gap, $\mathbf{P}_1$ and $\mathbf{P}_2$, define the
\emph{chord}
\begin{equation}
  \hat{\mathbf{e}} \;=\; \frac{\mathbf{P}_2-\mathbf{P}_1}
                              {\lVert\mathbf{P}_2-\mathbf{P}_1\rVert},
  \qquad
  L \;=\; \lVert\mathbf{P}_2-\mathbf{P}_1\rVert .
  \label{eq:chord}
\end{equation}
Axial (in-line) displacement components are obtained by projection onto the
chord,
\begin{equation}
  u_\parallel(x_i, t) \;=\; \mathbf{u}(x_i,t)\cdot\hat{\mathbf{e}},
  \label{eq:projection}
\end{equation}
which is the displacement component a DAS fibre would sense. Projection onto
the chord (rather than taking the $x$-component) avoids mixing transverse
motion into the axial estimate when the cable is not perfectly aligned with
the coordinate axes.

\paragraph{Sag measurement.}
% ldv_analysis/geometry.py :: measure_sag (use_3d=True, use_parabola=True)
For every scan point in the gap, the perpendicular distance from the chord
line is computed in full 3-D:
\begin{equation}
  s_i = (\mathbf{P}_i - \mathbf{P}_1)\cdot\hat{\mathbf{e}},
  \qquad
  d_i = \bigl\lVert (\mathbf{P}_i-\mathbf{P}_1) - s_i\,\hat{\mathbf{e}}
        \bigr\rVert ,
  \label{eq:perpdist}
\end{equation}
where $s_i$ is the arc coordinate along the chord and $d_i$ the perpendicular
offset. Since the endpoints lie on the chord by construction
($d(0) = d(L) = 0$), a parabola \emph{constrained through both endpoints} is
fitted,
\begin{equation}
  d(s) \;\approx\; -a\, s\,(s - L),
  \qquad
  a \;=\; \frac{\sum_i \phi_i\, d_i}{\sum_i \phi_i^2},
  \quad \phi_i = s_i (s_i - L),
  \label{eq:parabolafit}
\end{equation}
a one-parameter linear least-squares problem. The midpoint sag is the vertex
value of the fitted parabola,
\begin{equation}
  w_0 \;=\; d(L/2) \;=\; \frac{a L^2}{4} \Big|_{a<0} ,
  \label{eq:sagvalue}
\end{equation}
and the root-mean-square residual of the fit,
$\sigma_{w_0} = \sqrt{\tfrac{1}{N}\sum_i (d_i - d(s_i))^2}$, is retained as
the sag measurement uncertainty (it captures outlier scan points and
deviations of the true shape from the parabolic approximation). The parabola
is the small-sag shape of a beam/cable under uniform load, so this fit is
consistent with the analytical model of
Section~\ref{sec:proc:geometry:theta}. If the fit yields $a \ge 0$ (no
physical downward arch), the maximum raw perpendicular distance is used
instead.

\paragraph{The coupling parameter $\Theta$.}
\label{sec:proc:geometry:theta}
% ldv_analysis/config.py :: theta_from_sag, theta_from_material
The bending-stress-relief model \citep{ProbstSimone} predicts that the
static strain-transfer efficiency of a suspended segment is
\begin{equation}
  \eta_\mathrm{stat} \;=\; \frac{\delta x_\ell}{\delta L}
                    \;=\; \frac{1}{1+\Theta},
  \qquad
  \Theta \;=\; \frac{EA}{L}\,\frac{(C_1 x_{f0})^2}{E I C_2},
  \label{eq:etatheta}
\end{equation}
where $x_{f0}$ is the gravity-induced midpoint sag,
$C_1 = \int_0^L (\phi')^2\,\mathrm{d}x$ and
$C_2 = \int_0^L (\phi'')^2\,\mathrm{d}x$ are the mode-shape constants of the
first symmetric clamped--clamped deflection shape
$\phi(x) = \tfrac{1}{2}\bigl[1-\cos(2\pi x/L)\bigr]$. Substituting the
tension-free equilibrium sag
\begin{equation}
  x_{f0} \;=\; w_0 \;=\; \frac{\rho A g L^4}{4\pi^4 E I}
  \label{eq:sagtheory}
\end{equation}
into Eq.~\eqref{eq:etatheta} gives the \emph{material form}
\begin{equation}
  \Theta_\mathrm{mat} \;=\; \frac{\rho^2 g^2 A^3 L^8}{128\,\pi^8 E^2 I^3},
  \label{eq:thetamaterial}
\end{equation}
while keeping the \emph{measured} sag $w_0$ gives the \emph{sag form}
\begin{equation}
  \Theta_\mathrm{sag} \;=\; \frac{A\, w_0^2}{8 I}
  \;\overset{\text{circular}}{=}\; \frac{1}{2}\left(\frac{w_0}{r}\right)^2 ,
  \label{eq:thetasag}
\end{equation}
where the last equality holds for a circular cross-section with
$A = \pi r^2$ and $I = \pi r^4/4$. For the rectangular ribbon cable
(width $w$, height $h$, bending about the weak axis) $A = wh$ and
$I = w h^3/12$ are used, so Eq.~\eqref{eq:thetasag} correctly generalises
beyond circular cables. The sag form is preferred throughout this work
because it absorbs the actual (pretension-affected) deployment geometry; the
material form is reported for comparison. Young's modulus enters as the mean
of three independent tensile tests per cable; $\Theta_\mathrm{mat}$ is
evaluated once per test value and reported as mean~$\pm$~standard deviation.

\subsection{Strain and Elongation Estimation}
\label{sec:proc:strain}

The coupling efficiency compares the \emph{axial elongation of the cable}
with the \emph{elongation imposed at its boundaries}. Three independent
estimators of the cable elongation are implemented --- M1 per-segment 3-D
arc-length, M2 spatial gradient and M3 Fourier-domain gradient --- and their
agreement is a built-in consistency check of the processing.

\subsubsection{Method 1: Per-Segment 3-D Arc-Length Elongation}
\label{sec:proc:strain:m1}
% ldv_analysis/strain.py :: strain_3d_arclength
Method 1 estimates the elongation of each segment between two adjacent scan
points directly from their 3-D displacement vectors, without projecting the
whole field onto the global chord first. Let
$\mathbf{r}_i = \mathbf{P}_{i+1} - \mathbf{P}_i$ be the static segment
vector, $L_{0,i} = \lVert\mathbf{r}_i\rVert$ its rest length and
$\hat{\mathbf{e}}_i = \mathbf{r}_i / L_{0,i}$ its \emph{local} unit vector.
The deformed segment length is
\begin{equation}
  L_i(t) \;=\; \bigl\lVert \mathbf{r}_i
      + \Delta\mathbf{u}_i(t) \bigr\rVert ,
  \qquad
  \Delta\mathbf{u}_i(t) = \mathbf{u}_{i+1}(t) - \mathbf{u}_i(t).
\end{equation}
Because the dynamic displacements ($\sim\si{\micro\meter}$) are many orders
of magnitude smaller than the segment lengths ($\sim\si{\milli\meter}$), the
square root is linearised by a first-order Taylor expansion:
\begin{equation}
  L_i(t) - L_{0,i}
  \;=\; L_{0,i}\sqrt{1
        + \frac{2\,\mathbf{r}_i\!\cdot\!\Delta\mathbf{u}_i}{L_{0,i}^2}
        + \frac{\lVert\Delta\mathbf{u}_i\rVert^2}{L_{0,i}^2}} - L_{0,i}
  \;\approx\;
  \hat{\mathbf{e}}_i \cdot \Delta\mathbf{u}_i(t)
  \;\equiv\; \delta d_i(t).
  \label{eq:strainM1}
\end{equation}
Geometrically: \emph{the relative displacement of the two points is projected
onto the local segment direction} $\hat{\mathbf{e}}_i$ — only the component
of relative motion along the segment changes its length to first order; the
perpendicular component merely rotates it (a second-order length effect that
the linearisation deliberately drops, consistent with the small-strain
regime). For a sagged cable the $\hat{\mathbf{e}}_i$ follow the local cable
direction, so Method 1 measures the \emph{arc-length} change, which is
exactly the quantity a strained optical fibre integrates. The local strain
estimate of segment $i$ is $\varepsilon_i = \delta d_i / L_{0,i}$, and the
total cable elongation is the telescoping sum over the gap,
\begin{equation}
  \delta x_\ell(t) \;=\; \sum_{i\,\in\,\mathrm{gap}} \delta d_i(t).
  \label{eq:deltaxl}
\end{equation}
Method 1 is the primary estimator in this work: it makes no assumption
about the smoothness of the displacement field between scan points and is
insensitive to the chord-projection error for sagged geometries.

\subsubsection{Method 2: Spatial Gradient of the Axial Displacement}
% ldv_analysis/strain.py :: strain_spatial_gradient
Axial strain is the spatial derivative of the axial displacement along the
cable,
\begin{equation}
  \varepsilon(s,t) \;=\; \frac{\partial\, u_\parallel(s,t)}{\partial s},
  \label{eq:strainM2def}
\end{equation}
evaluated with second-order central finite differences on the non-uniform
scan-point grid. For interior point $i$ with spacings
$\Delta s_- = s_i - s_{i-1}$ and $\Delta s_+ = s_{i+1} - s_i$:
\begin{equation}
  \varepsilon_i \;=\;
    u_{\parallel,i+1}\,\frac{\Delta s_-}{\Delta s_+(\Delta s_+ + \Delta s_-)}
  + u_{\parallel,i}\,\frac{\Delta s_+ - \Delta s_-}{\Delta s_+ \Delta s_-}
  - u_{\parallel,i-1}\,\frac{\Delta s_+}{\Delta s_-(\Delta s_+ + \Delta s_-)} .
  \label{eq:strainM2}
\end{equation}
This reduces to the familiar $(u_{i+1}-u_{i-1})/(2\Delta s)$ for uniform
spacing and is exact for locally quadratic displacement fields. The total
elongation follows by trapezoidal integration over the gap,
$\delta x_\ell^{(2)}(t) = \int_\mathrm{gap} \varepsilon(s,t)\,\mathrm{d}s$.

\subsubsection{Method 3: Fourier-Domain Spatial Gradient}
% ldv_analysis/strain.py :: strain_fourier_gradient
As a cross-check with different numerical error characteristics, the axial
displacement is resampled onto a uniform grid $s_k$ (cubic interpolation)
and differentiated spectrally using the Fourier differentiation theorem,
\begin{equation}
  \varepsilon(s,t) \;=\;
  \mathcal{F}_s^{-1}\bigl\{\, i k\, \mathcal{F}_s\{u_\parallel\}\,\bigr\},
  \label{eq:strainM3}
\end{equation}
with $k$ the spatial angular wavenumber. Spectral differentiation is exact
for band-limited periodic signals; since the cable displacement is not
periodic over the gap, Gibbs-type artefacts can appear near the endpoints,
which is why only gap-interior stations enter the elongation integral. In
practice Method 3 agrees with Methods 1--2 in the smooth interior and serves
to expose discretisation artefacts of the finite-difference estimate.

\subsubsection{Reference Elongations}
% ldv_analysis/freqdomain.py :: build_coupling_signals (x_source)
Three references for the imposed boundary elongation $\delta L(t)$ are
available:
\begin{itemize}
  \item \textbf{Endpoint pair} (default): the span change measured on the
    cable itself,
    $\delta L = u_\parallel(\mathbf{P}_2) - u_\parallel(\mathbf{P}_1)$.
    It shares its endpoints with $\delta x_\ell$, so the efficiency isolates
    bending stress relief and is immune to cross-instrument calibration.
  \item \textbf{Shaker}: the chord-projected shaker-head displacement
    $\delta L = \mathbf{u}_\mathrm{sh}\cdot\hat{\mathbf{e}}$ (sign-corrected
    so that positive means stretching), i.e.\ the imposed ground motion.
  \item \textbf{Shaker-to-end}: the elongation of the sub-span between the
    shaker and the adjacent endpoint sensor, used to diagnose losses at the
    first contact.
\end{itemize}

\subsection{Time-Domain Coupling Efficiency (Quality Control)}
\label{sec:proc:timedomain}
% ldv_analysis/analysis.py :: per_frequency_qc

For inspection at a single excitation frequency $f_0$, the record is windowed
to $t_c(f_0) \pm n_\mathrm{cyc}/f_0$ (Eq.~\ref{eq:sweeptime},
$n_\mathrm{cyc}=5$), band-passed (Section~\ref{sec:proc:preprocessing}),
integrated to displacement, and the instantaneous efficiency is formed:
\begin{equation}
  \eta(t) \;=\; \frac{\delta x_\ell(t)}{\delta L(t)},
  \qquad
  \eta_\mathrm{med} = \operatorname{median}_t\, \eta(t)
  \quad\text{over } |\delta L(t)| > 0.05\,\max|\delta L| .
  \label{eq:etatime}
\end{equation}
The amplitude gate excludes the zero crossings of the reference, where the
ratio is ill-conditioned; the median suppresses the remaining outliers. A
phase-insensitive variant replaces both signals by their analytic-signal
envelopes (Section~\ref{sec:proc:amplitude}),
$\eta_\mathrm{env}(t) = |\mathcal{H}\{\delta x_\ell\}| /
|\mathcal{H}\{\delta L\}|$: comparing $\eta_\mathrm{med}$ and
$\eta_\mathrm{env}$ distinguishes whether a low efficiency is caused by an
amplitude deficit or by a phase lag between cable and reference.

These time-domain values are used for quality control and for illustrating
the mechanism at single frequencies; all summary efficiencies in the results
chapter are computed in the frequency domain
(Section~\ref{sec:proc:freqdomain}), which uses the whole record at once and
is considerably more stable.

\subsection{Frequency-Domain Transfer Functions}
\label{sec:proc:freqdomain}
% ldv_analysis/freqdomain.py

\subsubsection{Signals}
From the full (unfiltered) sweep record, integrated to displacement, three
signals are built:
\begin{align}
  X(t) &= \delta L(t) && \text{input: boundary span change,} \notag\\
  Y_1(t) &= \delta x_\ell(t) && \text{output 1: cable elongation
                                (Eq.~\ref{eq:deltaxl}),} \notag\\
  Y_2(t) &= u_{\perp,\mathrm{mid}}(t) && \text{output 2: midpoint deflection
                                along the static sag direction.} \notag
\end{align}
$H_\eta(f) = Y_1/X$ is the dynamic strain-transfer efficiency;
$H_\mathrm{mid}(f) = Y_2/X$ measures the intensity of the bending mode, whose
resonance structure explains the frequency dependence of $H_\eta$
\citep{ProbstSimone}.

\subsubsection{Welch Spectral Estimation}
% ldv_analysis/freqdomain.py :: welch_spectra
The record is divided into $K$ segments of length
$n_\mathrm{seg} = 2 f_s = 10\,000$ samples (\SI{2}{s}, frequency resolution
$\Delta f = \SI{0.5}{Hz}$) with 50\,\% overlap, each tapered with a Hann
window $w[n]$. With $X_k(f)$, $Y_k(f)$ the windowed DFTs of segment $k$,
the auto- and cross-spectral densities are the segment averages
\begin{equation}
  S_{XX}(f) = \frac{1}{K}\sum_{k=1}^{K} |X_k(f)|^2 , \qquad
  S_{YY}(f) = \frac{1}{K}\sum_{k=1}^{K} |Y_k(f)|^2 , \qquad
  S_{XY}(f) = \frac{1}{K}\sum_{k=1}^{K} X_k^*(f)\, Y_k(f),
  \label{eq:welch}
\end{equation}
(up to the standard window normalisation, which cancels in every ratio used
below). The averaging over $K \approx 11$ segments reduces the variance of
each spectral estimate by $\approx 1/K$ — this is what makes the Welch-based
transfer functions far smoother than any single-FFT estimate. The
\emph{ordinary coherence}
\begin{equation}
  \gamma^2(f) \;=\; \frac{|S_{XY}(f)|^2}{S_{XX}(f)\, S_{YY}(f)}
  \;\in\; [0, 1]
  \label{eq:coherence}
\end{equation}
measures the fraction of output power linearly explained by the input;
$\gamma^2 < 1$ indicates noise, nonlinearity, or unmeasured inputs.

\subsubsection{FRF Estimators: Direct, $H_1$, $H_2$, $H_v$}
\label{sec:proc:estimators}
% ldv_analysis/freqdomain.py :: direct_transfer, estimate_H

\paragraph{Direct estimate.}
The conceptually simplest frequency response estimate is the ratio of the
full-record windowed Fourier transforms,
\begin{equation}
  H_\mathrm{direct}(f) \;=\;
  \frac{\mathcal{F}\{w\,Y\}(f)}{\mathcal{F}\{w\,X\}(f)} ,
  \label{eq:Hdirect}
\end{equation}
evaluated on the dense grid $\Delta f_\mathrm{d} = f_s/N = \SI{0.083}{Hz}$.
Every frequency bin is a single realisation: the estimate is unbiased in the
noise-free limit but has no variance reduction, so amplitude and phase
scatter strongly wherever excitation or response is weak.

Comparing Eq.~\eqref{eq:Hdirect} directly against $H_1$ would conflate two
distinct effects — the estimator algebra and the amount of spectral averaging.
To separate them, the direct spectra are smoothed along the \emph{frequency}
axis (Daniell smoothing) rather than over time segments: with a normalised
kernel $k$ of $m$ bins,
\begin{equation}
  \tilde S_{AB}(f_i) = \sum_{j} k_j\, A^*(f_{i+j})\, B(f_{i+j}),
  \qquad
  H_\mathrm{direct}^{\mathrm{sm}}(f) = \frac{\tilde S_{XY}(f)}{\tilde S_{XX}(f)} .
  \label{eq:Hdirectsmooth}
\end{equation}
Smoothing the products \emph{before} the division makes the estimate
input-power weighted, so bins with weak $|X|$ contribute little instead of
diverging; this also removes the positive bias that the unweighted ratio
carries at low SNR. A kernel with $\sum_j k_j = 1$ averages
$n_\mathrm{eff} = 1/\sum_j k_j^2$ independent bins, which sets both the
variance reduction and the effective resolution bandwidth. The kernel width is
therefore chosen so that $n_\mathrm{eff}\,\Delta f_\mathrm{d}$ equals the
\emph{equivalent noise bandwidth} of the Welch estimate,
\begin{equation}
  B_\mathrm{ENBW} = f_s \frac{\sum_n w[n]^2}{\left(\sum_n w[n]\right)^2}
  = 1.5\,\Delta f = \SI{0.75}{Hz}
  \quad\text{(Hann)},
  \label{eq:enbw}
\end{equation}
i.e.\ $n_\mathrm{eff} \approx 9$ dense bins. Note that
$B_\mathrm{ENBW} > \Delta f$: the Hann taper spreads each spectral line over
1.5 bins, so a Welch estimate does not resolve to its nominal bin spacing, and
matching $B_\mathrm{ENBW}$ rather than $\Delta f$ is what puts the two
estimates on an equal footing. With this choice the smoothed direct estimate
and $H_1$ agree to within a few percent in $|H|$ while retaining the dense
frequency grid; the residual difference is attributable to the estimator
algebra alone. Both the raw and the smoothed direct curves are shown, so the
figure separates the effect of averaging (raw $\to$ smoothed) from the effect
of the estimator (smoothed $\to$ $H_1$).

\paragraph{$H_1$ — noise on the output.}
Assume the measured output contains additive noise uncorrelated with the
input, $Y = H X + N$. Minimising the output residual over the segment
ensemble,
$\min_H \sum_k |Y_k - H X_k|^2$, gives the least-squares solution
\begin{equation}
  H_1(f) \;=\; \frac{S_{XY}(f)}{S_{XX}(f)} .
  \label{eq:H1}
\end{equation}
Since $\mathbb{E}\{X^* N\} = 0$, the cross-spectrum averages the output noise
away and $H_1$ is unbiased under output noise. If the \emph{input} channel is
noisy, $S_{XX}$ is inflated and $|H_1|$ \emph{under}estimates the true FRF.

\paragraph{$H_2$ — noise on the input.}
Assume instead noise on the input, $X = X_0 + M$ with $Y = H X_0$.
Minimising the input-referenced residual
$\min_H \sum_k |X_k - Y_k / H|^2$ yields
\begin{equation}
  H_2(f) \;=\; \frac{S_{YY}(f)}{S_{YX}(f)}
        \;=\; \frac{S_{YY}(f)}{S_{XY}^*(f)} .
  \label{eq:H2}
\end{equation}
$H_2$ is unbiased under input noise but \emph{over}estimates $|H|$ when the
output is noisy. The two estimators bracket the true FRF,
\begin{equation}
  |H_1(f)| \;\le\; |H(f)| \;\le\; |H_2(f)|,
  \qquad
  \frac{|H_1(f)|}{|H_2(f)|} = \gamma^2(f),
  \label{eq:bracket}
\end{equation}
so their separation is itself a coherence display: the estimators agree
exactly where $\gamma^2 = 1$ and diverge where the data are noisy.

\paragraph{$H_v$ — noise on both channels (total least squares).}
When both channels carry comparable noise — realistic here, since input and
output are both derived from LDV velocities through the same processing —
the total-least-squares estimator distributes the residual over both
channels. Writing the spectral sample covariance of the vector $(X, Y)$ per
frequency bin,
\begin{equation}
  \mathbf{C}(f) \;=\;
  \begin{pmatrix} S_{XX} & S_{XY}^* \\ S_{XY} & S_{YY} \end{pmatrix},
\end{equation}
the TLS solution is the direction of the eigenvector belonging to the
\emph{smallest} eigenvalue $\lambda_-$ of $\mathbf{C}$ (the noise
eigenvalue). For the $2\times 2$ Hermitian matrix this eigenvalue is
\begin{equation}
  \lambda_-(f) \;=\; \frac{S_{XX}+S_{YY}}{2}
    - \sqrt{\left(\frac{S_{XX}-S_{YY}}{2}\right)^{\!2} + |S_{XY}|^2},
\end{equation}
and solving the eigenvector relation for the FRF gives the closed form
\begin{equation}
  H_v(f) \;=\; \frac{S_{YY}(f) - \lambda_-(f)}{S_{XY}^*(f)} .
  \label{eq:Hv}
\end{equation}
One verifies directly that in the noise-free case ($\lambda_- = 0$,
$S_{YY} = |H|^2 S_{XX}$, $S_{XY} = H S_{XX}$) Eq.~\eqref{eq:Hv} returns
exactly $H$, and that $|H_1| \le |H_v| \le |H_2|$ holds bin-by-bin. All
three Welch estimators share the same phase,
$\arg H_1 = \arg H_2 = \arg H_v = \arg S_{XY}$, so the estimator choice
affects only the magnitude.

\paragraph{Phase presentation.}
Phases are reported \emph{wrapped} to $(-180^\circ, 180^\circ]$. Unwrapped
phase curves accumulate spurious $\pm 360^\circ$ jumps whenever
low-coherence bins randomise the phase, which makes them diverge without
physical meaning; the wrapped display is immune to this and the physically
relevant transition (the $\approx 180^\circ$ flip across the bending
resonance) remains fully visible.

\subsubsection{Band-Averaged Efficiency and Coherence Gating}
% ldv_analysis/freqdomain.py :: band_mean_eta
The quasi-static efficiency reported per configuration is the coherence-gated
band average
\begin{equation}
  \bar\eta \;=\; \bigl\langle\, |H_\eta(f)| \,\bigr\rangle_{f \in B},
  \qquad
  B = \{ f \in [f_\mathrm{min}, f_\mathrm{max}] :
         \gamma^2_\eta(f) \ge 0.7 \},
  \label{eq:bandeta}
\end{equation}
with $[f_\mathrm{min}, f_\mathrm{max}] = [5, 25]\,\si{Hz}$ well below the
first resonance. The gate removes noise-dominated bins that would otherwise
bias the average; the standard deviation over $B$ is reported as the error
bar.

\subsubsection{Single-Frequency Cross-Check and Resonance Characterisation}
% ldv_analysis/freqdomain.py :: stft_point_transfer, fit_lorentzian
As an independent validation, point estimates
$H(f_0) = \mathcal{F}\{w\,Y\}(f_0)\,/\,\mathcal{F}\{w\,X\}(f_0)$ are computed
from short Hann-tapered windows centred on the sweep arrival $t_c(f_0)$ and
overlaid on the Welch curves.

The fundamental resonance $f_1$ is identified as the global maximum of
$|H_\mathrm{mid}(f)|$; a driven single-degree-of-freedom magnitude response
\begin{equation}
  |H(f)| \;=\; \frac{G\, f_1^2}
   {\sqrt{\left(f_1^2 - f^2\right)^2 + \left(f_1 f / Q\right)^2}}
  \label{eq:sdof}
\end{equation}
is fitted in a window around the peak, yielding the quality factor $Q$ and
damping ratio $\zeta = 1/(2Q)$. The theoretical clamped--clamped fundamental
\begin{equation}
  f_1^\mathrm{th} \;=\; \frac{\alpha_1^2}{2\pi L^2}\sqrt{\frac{EI}{\rho A}},
  \qquad \alpha_1^2 = 4.730^2 \approx 22.37,
  \label{eq:f1theory}
\end{equation}
is quoted for comparison (cf.~Section~\ref{sec:proc:modal}).

\subsection{Amplitude Extraction}
\label{sec:proc:amplitude}
% ldv_analysis/amplitude.py

Beyond ratios, the absolute displacement and strain amplitudes are needed as
reference values for the numerical simulations. Two extraction routes are
implemented and cross-validated.

\paragraph{Time domain: analytic-signal envelope.}
For a windowed, band-passed trace $x(t)$ the analytic signal is
$x_a(t) = x(t) + i\,\mathcal{H}\{x\}(t)$ with $\mathcal{H}$ the Hilbert
transform; its magnitude $|x_a(t)|$ is the instantaneous envelope. The
representative amplitude is the envelope median (robust against the
filter-edge droop at the window boundaries); envelope maximum and
$\sqrt{2}\,\mathrm{RMS}$ variants are available for comparison.

\paragraph{Frequency domain: band-energy read-out.}
The velocities are integrated once per dataset; at each target frequency
$f_0$ the displacement signal is windowed around $t_c(f_0)$, Hann-tapered
and Fourier transformed. Because the log-sweep smears the tone over several
bins (especially at low frequency, where the window spans many seconds), the
amplitude is recovered from the \emph{energy in the analysis band}
$B = [f_0 - 5, f_0 + 5]\,\si{Hz}$ rather than from a single bin. For a tone
of peak amplitude $A$ under a window $w[n]$, Parseval's theorem gives
\begin{equation}
  A \;=\; 2\sqrt{\frac{\sum_{k \in B} |X_k|^2}{N \sum_n w[n]^2}} ,
  \label{eq:bandenergy}
\end{equation}
which is exact for an isolated in-band tone. This route avoids the
per-frequency band-pass filtering of the full record and is roughly two
orders of magnitude faster, enabling dense frequency grids; it agrees with
the envelope route to within a few percent below $\approx\SI{100}{Hz}$.

\subsection{Extension--Compression Asymmetry}
\label{sec:proc:asymmetry}
% ldv_analysis/asymmetry.py

A sagged segment stiffens when its endpoints are pulled apart (the sag is
straightened) and softens when they are pushed together (the sag deepens);
numerical simulations show correspondingly different responses for the two
deformation directions. A perfectly linear system responds
mirror-symmetrically to the two half-cycles of a sinusoidal drive, so any
measurable asymmetry is a signature of this geometric nonlinearity. Three
complementary detectors are evaluated per excitation frequency:

\begin{enumerate}
  \item \textbf{Half-cycle conditional efficiency.} The samples of the
  reference elongation are split by sign — extension
  ($\delta L > 0$, endpoints separating) versus compression
  ($\delta L < 0$) — with a $25\,\%$-of-maximum amplitude gate excluding the
  zero crossings. The efficiencies
  \begin{equation}
    \eta_\mathrm{ext} = \operatorname*{median}_{\delta L(t) > \delta_g}
      \frac{\delta x_\ell(t)}{\delta L(t)},
    \qquad
    \eta_\mathrm{comp} = \operatorname*{median}_{\delta L(t) < -\delta_g}
      \frac{\delta x_\ell(t)}{\delta L(t)}
    \label{eq:etahalfcycle}
  \end{equation}
  coincide for a linear system; $\eta_\mathrm{ext} \ne \eta_\mathrm{comp}$
  quantifies the directional coupling difference.

  \item \textbf{Peak asymmetry.} From the mean positive-peak amplitude
  $A_+$ and mean negative-peak magnitude $A_-$ of the cable elongation
  $\delta x_\ell(t)$:
  \begin{equation}
    \alpha \;=\; \frac{A_+ - A_-}{A_+ + A_-} \;\in\; [-1, 1],
    \label{eq:peakasym}
  \end{equation}
  with $\alpha > 0$ meaning the cable elongates further during extension
  half-cycles than it shortens during compression.

  \item \textbf{Even-harmonic distortion.} Expanding the response in powers
  of the drive, $y = c_1 x + c_2 x^2 + c_3 x^3 + \dots$ with
  $x = A\cos\omega t$, the quadratic term produces
  $c_2 A^2 \cos^2\omega t = \tfrac{c_2 A^2}{2}(1 + \cos 2\omega t)$ — energy
  at \emph{twice} the drive frequency — whereas the cubic (symmetric) term
  appears at $3\omega$. Since a sign-asymmetric stiffness is exactly a
  quadratic term, the ratio
  \begin{equation}
    R_2 \;=\; \frac{|Y(2 f_0)|}{|Y(f_0)|}
    \label{eq:h2ratio}
  \end{equation}
  isolates the asymmetric part of the response independently of the
  reference channel; $R_3 = |Y(3f_0)|/|Y(f_0)|$ is reported alongside to
  attribute nonlinearity to the symmetric (softening/stiffening) branch.
\end{enumerate}

For the linearity data set (Section~\ref{sec:proc:linearity}) the half-cycle
efficiencies are additionally evaluated in bins of the instantaneous drive
amplitude, revealing whether an asymmetry develops progressively with
amplitude (e.g.\ slack forming on the compression half-cycle only).

\subsection{Modal Analysis}
\label{sec:proc:modal}
% ldv_analysis/modal.py

\subsubsection{Detection Fingerprint and Spectral Smoothing}
% modal.spectral_fingerprint
Transverse resonances are detected on a \emph{fingerprint}: per-sensor
spectra of the transverse velocities $v_y, v_z$ over the gap, collapsed
across sensors into the total-power detection signal
\begin{equation}
  P(f) \;=\; \sum_{s\,\in\,\mathrm{gap}}
      \Bigl( |V_y(f; s)|^2 + |V_z(f; s)|^2 \Bigr).
  \label{eq:fingerprint}
\end{equation}
The per-sensor spectra are Welch estimates
($n_\mathrm{seg} = 2 f_s$, $\Delta f = \SI{0.5}{Hz}$, Hann, 50\,\% overlap,
$K \approx 11$ averages) rather than raw single-record FFTs: the
$\mathrm{var} \propto 1/K$ averaging suppresses the bin-to-bin fluctuations
that previously caused the peak picker to fire on noise wiggles. An optional
moving-average (boxcar) smoothing over a configurable bandwidth can be
applied on top. Resonances are then the peaks of $P(f)$ (normalised to its
in-band maximum) exceeding a prominence threshold with a minimum mutual
spacing (\texttt{scipy.signal.find\_peaks}).

\paragraph{Input-bias correction.}
% modal.shaker_input_spectrum, normalize_by_input
The shaker head does not excite purely axially: its own recording shows
frequency-dependent energy in the transverse components. A spectral peak in
the cable response may therefore reflect a peak in the \emph{excitation}
rather than a cable resonance. Two diagnostics address this. First, the
input directionality is quantified by the per-frequency power fractions of
the shaker-head components,
\begin{equation}
  \phi_c(f) \;=\; \frac{|V_c^\mathrm{sh}(f)|^2}
   {\sum_{c' \in \{x,y,z\}} |V_{c'}^\mathrm{sh}(f)|^2},
  \qquad c \in \{x, y, z\}.
  \label{eq:directionality}
\end{equation}
Second, the detection signal can be normalised by the total shaker input
power, $P_\mathrm{norm}(f) = P(f) / P_\mathrm{sh}(f)$ — an FRF-like quantity
in which input coloration cancels: peaks that survive the normalisation are
genuine output-per-input amplifications (resonances), while peaks that
disappear were inherited from the excitation spectrum.

\subsubsection{Complex Mode-Shape Extraction}
% modal.extract_mode_shape
For each detected resonance $f_n$, the transverse velocities are band-passed
to $f_n(1 \pm 0.1)$, integrated to displacement, and the complex Fourier
coefficient at the bin nearest $f_n$ is read per sensor:
\begin{equation}
  \varphi_y(s) = U_y(f_n; s), \qquad \varphi_z(s) = U_z(f_n; s).
  \label{eq:modeshape}
\end{equation}
The complex-valued shapes retain the relative phase between sensors, so
travelling-wave content and sign flips across nodes are preserved.

\subsubsection{Euler--Bernoulli Reference Modes}
% modal.clamped_clamped_mode
The gap is modelled as a doubly-clamped uniform beam. Free vibration of an
Euler--Bernoulli beam obeys $EI\, w'''' = \rho A\, \omega^2 w$; with clamped
boundary conditions $w(0)=w'(0)=w(L)=w'(L)=0$, non-trivial solutions exist
only at the roots $\alpha_n$ of the frequency equation
\begin{equation}
  \cos\alpha \cosh\alpha \;=\; 1,
  \qquad
  \alpha_n = 4.730,\ 7.853,\ 10.996,\ 14.137,\ 17.279,\ \dots
  \label{eq:freqequation}
\end{equation}
with natural frequencies and mode shapes
\begin{equation}
  f_n \;=\; \frac{\alpha_n^2}{2\pi L^2} \sqrt{\frac{EI}{\rho A}},
  \qquad
  \varphi_n(\xi) = \bigl[\cosh\alpha_n\xi - \cos\alpha_n\xi\bigr]
   - \sigma_n \bigl[\sinh\alpha_n\xi - \sin\alpha_n\xi\bigr],
  \label{eq:ebmodes}
\end{equation}
where $\xi = x/L \in [0,1]$ and
$\sigma_n = (\cosh\alpha_n - \cos\alpha_n)/(\sinh\alpha_n - \sin\alpha_n)$.
The frequency ratios $f_n / f_1 = (\alpha_n/\alpha_1)^2 =
1,\ 2.76,\ 5.40,\ 8.93,\ \dots$ provide a material-free consistency check of
the mode ordering.

\subsubsection{Modal Assurance Criterion}
% modal.mac
Each measured shape is assigned to a theoretical mode by the Modal Assurance
Criterion,
\begin{equation}
  \operatorname{MAC}(\boldsymbol{\varphi}_m, \boldsymbol{\varphi}_t)
  \;=\;
  \frac{\bigl|\boldsymbol{\varphi}_m^{\mathsf{H}}
              \boldsymbol{\varphi}_t\bigr|^2}
       {\bigl(\boldsymbol{\varphi}_m^{\mathsf{H}}\boldsymbol{\varphi}_m\bigr)
        \bigl(\boldsymbol{\varphi}_t^{\mathsf{H}}\boldsymbol{\varphi}_t\bigr)}
  \;\in\; [0, 1],
  \label{eq:mac}
\end{equation}
where $\boldsymbol{\varphi}_m$ is the measured complex shape (the dominant
polarization component, sampled at the sensor positions),
$\boldsymbol{\varphi}_t$ the theoretical shape $\varphi_n(\xi_i)$ evaluated
at the same positions, and $^{\mathsf{H}}$ the conjugate transpose. The MAC
is the squared modulus of the normalised inner product — a
\emph{shape-correlation} coefficient: it equals 1 when the vectors are
collinear (identical shape up to any complex scale), and 0 when they are
orthogonal. Because a global complex factor cancels, the measured shape
needs no phase alignment or amplitude scaling before comparison, and a
$180^\circ$ phase flip across a node is handled correctly. The measured
resonance is assigned the mode order with the largest MAC; low best-MAC
values flag shapes that no beam mode explains (e.g.\ rigid-body-like motion
from mount compliance or shaker cross-talk).

\subsubsection{Polarization Analysis}
% modal.polarization_ellipse
At each sensor the pair $(\varphi_y, \varphi_z)$ defines an ellipse traced
by $\Re\{(\varphi_y, \varphi_z)\, e^{i\omega t}\}$ in the $y$--$z$ plane.
Decomposing into counter-rotating circular components
\begin{equation}
  A_+ = \tfrac{1}{2}\bigl(\varphi_y + i\,\varphi_z\bigr),
  \qquad
  A_- = \tfrac{1}{2}\bigl(\varphi_y^* + i\,\varphi_z^*\bigr)
  \quad\Longrightarrow\quad
  \text{major} = |A_+| + |A_-|,\quad
  \text{minor} = \bigl||A_+| - |A_-|\bigr|,
  \label{eq:ellipse}
\end{equation}
with axis ratio $\mathrm{minor}/\mathrm{major}$ (0 = linear, 1 = circular),
orientation $\theta = \tfrac{1}{2}(\arg A_+ + \arg A_-)$ and handedness
$\operatorname{sign}\Im(\varphi_y \varphi_z^*)$. A resonance is classified
linear / elliptical / circular by the amplitude-weighted majority over the
significant sensors.

\subsection{Linearity Analysis}
\label{sec:proc:linearity}
% ldv_analysis/analysis.py :: linearity_qc

A dedicated test signal probes amplitude dependence: eight discrete
frequencies (5--400\,\si{Hz}), each driven for \SI{3.5}{s} with a linear
amplitude ramp $0 \to 1$ followed by a \SI{0.5}{s} fade. Each segment is
extracted in full, band-passed ($f_0 \pm \SI{5}{Hz}$), integrated, and the
efficiency $\eta(t)$ (Eq.~\ref{eq:etatime}) is evaluated against the
\emph{instantaneous drive amplitude}, defined as the analytic-signal envelope
$a(t) = |\mathcal{H}\{\delta L\}(t)|$. Because the ramp sweeps amplitude
monotonically, plotting the binned median of $\eta$ against $a$ directly
tests linearity: a flat curve means amplitude-independent coupling, a slope
or hysteresis loop indicates nonlinearity (e.g.\ frictional slip at the
mounts above a threshold amplitude). A complementary slip indicator divides
the envelope of a single cable-sensor velocity by the envelope of the shaker
reference; for slip-free coupling this ratio stays constant along the ramp.
The frequency-domain cross-check estimates $H_1$
(Eq.~\ref{eq:H1}) on each mono-frequency segment and compares
$|H_1(f_0)|$ with the time-domain median efficiency.

\subsection{Uncertainty Treatment}
\label{sec:proc:uncertainty}

\begin{itemize}
  \item \textbf{Young's modulus:} three tensile tests per cable;
    $\Theta_\mathrm{mat}$ is evaluated per test and reported as mean $\pm$
    standard deviation. Since $\Theta_\mathrm{mat} \propto E^{-2}$, a
    relative modulus error $\epsilon_E$ propagates as
    $2\,\epsilon_E$ into $\Theta$.
  \item \textbf{Gap length:} the mounting-gap reading error is
    $\sigma_L = \SI{0.5}{cm}$. The material form is extremely sensitive,
    $\Theta_\mathrm{mat} \propto L^8 \Rightarrow
    \sigma_\Theta/\Theta = 8\,\sigma_L/L$ (40\,\% at $L=\SI{10}{cm}$!),
    which is the quantitative reason the sag-based
    $\Theta_\mathrm{sag}$ — independent of $L$ — is preferred.
  \item \textbf{Sag:} the RMS residual of the constrained parabola fit
    (Eq.~\ref{eq:parabolafit}) serves as $\sigma_{w_0}$; with
    $\Theta_\mathrm{sag} \propto w_0^2$,
    $\sigma_\Theta/\Theta = 2\,\sigma_{w_0}/w_0$.
  \item \textbf{Efficiency:} the standard deviation of $|H_\eta|$ over the
    coherence-gated quasi-static band (Eq.~\ref{eq:bandeta}); the
    $H_1$/$H_2$ bracket (Eq.~\ref{eq:bracket}) provides an additional
    systematic error envelope.
\end{itemize}

% ─── Bibliography hooks ─────────────────────────────────────────────────────
% \citep{ProbstSimone}  : S. Probst et al., "Unburied DAS on the Moon:
%                         Modeling Ground-to-Cable Strain Transfer" (draft).
% Probst et al. 2026    : Controlled Source DAS Coupling Tests, Earth and
%                         Space Science 13(3), e2025EA004817.
% Welch 1967            : P. Welch, IEEE Trans. Audio Electroacoust. 15(2).
% Allemang 2003         : R. Allemang, "The Modal Assurance Criterion --
%                         Twenty Years of Use and Abuse", Sound & Vibration.
% Bendat & Piersol 2010 : Random Data: Analysis and Measurement Procedures
%                         (H1/H2/Hv estimators, coherence).
